\chapter{Verbs}\label{ch:verbs}

\section{Introduction}

The verb is the morphological nucleus of the Iridian clause, and upon it falls
the principal burden of grammatical expression. Where nouns inflect for case and number, verbs
express the event's aspectual contour, the configuration of its participants,
and the speaker's modal orientation. Every finite verb bears the aspect-alignment suffix,
which is the sole obligatory element of the verbal complex; beyond this the verb
may carry markers of voice, mood, and evidentiality, in an order fixed by the
morphological template set out in the following section. Tense, gender, number
agreement, and pronominal cross-reference find no expression on the Iridian
verb.

\section{The morphological template}\label{sec:morph-template}

The structure of a finite verb in Iridian is given by the following template,
wherein parentheses enclose optional elements and the sole obligatory position
is that of the aspect-alignment suffix:

\medskip
\begin{center}
(\Neg{}) + \Stem{} + (\Antip{} / \Appl{}) + (\Caus{}) + (\textsc{orientation}) + \synhead{Asp.Align} + (\textsc{mood}) + (\textsc{evid})
\end{center}
\medskip

The \synhead{Asp.Align} slot fuses aspect and alignment into a single morpheme
and is the one element no finite verb may lack. The mood slot may be analysed as
containing a null morpheme in the indicative. The negative prefix \ird{za-}
stands before the entire stem complex. The indirect evidential has default
\ird{-al} after a consonant and contextual \ird{-l} after a vowel;
it follows finite aspect and mood and precedes the linker
(§\,\ref{sec:vm-evid}). Person is not morphologically expressed on the Iridian
verb: there is no pronominal clitic series and no agreement layer. Overt person
reference is conveyed by full pronoun DPs, demonstratives in the third person,
and zero anaphora where discourse recoverability licenses argument omission.

The \Stem{} position admits a derived lexical stem as well as a simple root.
The lexical preverbs \ird{de-}, \ird{mně-}, \ird{pře-},
\ird{le-}, \ird{ra-}, \ird{vra-}, \ird{ša-}, \ird{na-},
\ird{ru-}, and marginal \ird{ne-} belong inside that stem, as in
\ird{mněktař-} \trsl{read through} and \ird{lenočk-} \trsl{run out},
before the valency and orientation suffixes
(§\,\ref{sec:lexical-preverbs}). They modify the lexical event description
without supplying an Asp.Align value or changing the alignment domain;
indicative \ird{za-} attaches outside the complete derived stem.


\subsection{The Mirror Principle and functional projections}
\label{sec:vm-overview}

\noindent Baker's (1985) \term{Mirror Principle} predicts that morphological
derivations reflect the order of syntactic composition. The Iridian analysis
distinguishes argument-building heads from the higher orientation head realised
by \ird{-n-} (\ird{-an-} after a rhotic causative), so that
causativisation precedes the selection of the derived patient in \ird{piášt-r-an-ik}. Cumulative exponence and conditioned allomorphy
are specified with the individual constructions; the correspondence is not a
requirement that every head have a separate segment. Reading outward from the
root, the suffixal sequence has the following organisation. The separation
of event introduction, external argument introduction and case follows the
analytical distinctions surveyed by Harley \cite[§§1.3--1.4]{harley2017}; the
particular ordering and licensing conditions stated here belong to Iridian
and are summarised in Table~\ref{tab:vm-projections}.

\begin{table}
  \footnotesize
  \caption{Morphological template positions and corresponding
  syntactic projections (innermost to outermost).}\label{tab:vm-projections}
  \begin{tabularx}{\textwidth}{l l L l}
    \toprule
    {\sc position} & {\sc morpheme} & {\sc syntactic projection} &
    {\sc see} \\
    \midrule
    Root         & \Stem{}           & $\surd$P / \synhead{VP} (lexical predicate)
                 & §\,\ref{sec:vm-root} \\
    Inner valency & \Antip{} / \Appl{} & \synhead{VoiceP}$_{\mathrm{AP}}$ / \synhead{ApplP}
                 & §\,\ref{sec:vm-antipassive}, §\,\ref{sec:vm-n} \\
    Causative    & \Caus{}           & \synhead{CausP}
                 & §\,\ref{sec:vm-caus} \\
    Orientation & \ird{-n-}/\ird{-an-}, $\varnothing$ & \synhead{VoiceP}$_{\mathrm{orient}}$
                 & §\,\ref{sec:vm-n} \\
    Asp.Align    & \synhead{Asp.Align}.\textsc{align} & \synhead{AspP} (fused Asp + Case)
                 & §\,\ref{sec:vm-asp} \\
    Neg (prefix) & \Neg{}            & \synhead{NegP} (see §\,\ref{sec:vm-neg})
                 & §\,\ref{sec:vm-neg} \\
    Mood         & \textsc{mood}     & \synhead{MoodP} / \synhead{IP}
                 & §\,\ref{sec:vm-moodp} \\
    Evidential   & \ird{-al}/\ird{-l} & \synhead{EvidP}$_\text{inner}$
                 & §\,\ref{sec:vm-evid} \\
    \bottomrule
  \end{tabularx}
\end{table}

\noindent The negative marker \Neg{} precedes the root, although Neg is
analysed as structurally higher than Asp. Its preverbal position is attributed
to post-syntactic attachment to the left edge of the verbal complex
(§\,\ref{sec:vm-neg}).

\subsubsection{The root and VP}
\label{sec:vm-root}

For a simple verb, the lexical stem realises a root with a verbal categoriser
$v$. Categorisation, argument introduction and case licensing are distinct
operations in this analysis: the categoriser supplies verbal category, while
the predicate's lexical frame and any argument-building heads determine its
participants. An actor-introducing layer belongs to the inner predicate,
above Appl where present and below Caus; its ordinary exponent is zero.
The outer orientation head consequently need not introduce the actor whose
case its configuration determines.
A finite verb requires Asp.Align, and a bare stem cannot constitute an
independent finite predicate. Bound stems also occur in licensed non-finite
and nominal constructions, whose selecting heads need not include finite
Asp.Align (§\,\ref{sec:nominalised}).

A preverb--root combination is analysed here as a derived lexical predicate
within this innermost domain. Its internal prefixal order does not introduce
another member of the suffixal functional sequence in
Table~\ref{tab:vm-projections}; the ordinary Voice, causative, and aspectual
heads combine with the resulting stem. Prefixation preserves the base's
thematic class and construction-selected stem grades: \ird{nočk-},
\ird{stoj-}, and \ird{vad-} take ordinary progressive \ird{-ime} in
\ird{lenóčkime}, \ird{rustójime}, and \ird{navádime}, respectively.
No additional aspectual exponent or secondary imperfectivisation intervenes.

The \synhead{VP} formed by the root + its arguments constitutes the innermost
layer of the verbal clause. In a transitive clause the internal argument
(patient) is merged as the complement of V:

\begin{center}
\begin{forest}
  syntax tree
  [\synhead{VP}
    [{\synhead{DP}\\{\footnotesize patient}}]
    [V
      [{\textit{root}}]
    ]
  ]
\end{forest}
\end{center}

\noindent Because \synhead{VP} is head-final
(§\,\ref{sec:th-head-finality}), this complement surfaces to the left of the
verbal head at every embedded level. The preverbal position of the object is
therefore not derived by movement but is the base-generated configuration of the
head-final \synhead{VP}.



\section{The verb stem and thematic classes}

Iridian verbs admit no citation form. Lexicographers record the bare stem with a
following hyphen to indicate its character as a bound form requiring at least
one suffix; the verb rendered in English as \textit{to eat} is thus entered as
\ird{piašt-}, where \ird{t} constitutes the \term{thematic consonant}.

The thematic consonant, being the final consonant of the stem, is phonologically
active: it undergoes regular alternations when certain suffixes are adjoined.
Suffixes select distinct mutation classes: antipassive \ird{-aš-} triggers
\term{sibilant replacement}, whereas simultaneous \ird{-ěc} and conjunctive
\ird{-e} select the front grade; irrealis \ird{-ěv} selects the
yod grade (§\,\ref{sec:suffix-classes}). The grade is a property of the
construction rather than an automatic effect of the following vowel.
The five thematic classes in Table~\ref{tab:verbclasses} govern sibilant
replacement only; a listed stem shape does not independently establish that
the lexeme permits the antipassive argument frame.

\begin{table}
	\footnotesize
	\caption{Thematic consonant classes.}\label{tab:verbclasses}
	\begin{tabularx}{\textwidth}{l L L L}
		\toprule
		{\sc class} & {\sc thematic consonants} & {\sc change before \Antip{}} & {\sc example} \\
		\midrule
		I    & -t, -k, -c, -č  & → -č                    & \ird{piašt-} → \ird{piášč-} \\
		II   & -s, -š          & → -š                    & \ird{tras-} → \ird{tráš-} \\
		II-A & -l, -p          & add -š                  & \ird{dal-} → \ird{dalš-} \\
		III  & -d, -h, -j, -z, -ž & → -ž               & \ird{vad-} → \ird{váž-} \\
		III-A & all remaining  & add -ž                  & \ird{ščen-} → \ird{ščenž-} \\
		\bottomrule
	\end{tabularx}
\end{table}

Loan stems with surface \orth{g} are assigned analogically to the voiced
velar/fricative series and do not require a native phoneme /g/. Palatalising
suffixes do \emph{not} follow Table~\ref{tab:verbclasses}; they invoke the
consonantal-grade system described in
Section~\ref{sec:suffix-classes} of Chapter~\ref{ch:phon}.

Verbal stem grades are selected by the morphological construction and the lexical
class of the stem. Fleeting-vowel stems have a reduced grade, as in
\ird{razeň-}/\ird{razn-}; some modal constructions select a raised grade,
as in raised-strong \ird{zděk-}/\ird{zdíč-}; the perfective suffix has its own
\ird{-ek}/\ird{-ik-} relation. These changes preserve quantity unless an
independent quantity grade is specified. Consonant mutation applies to the stem
or aspect suffix selected in that domain, followed by syllabification and the
ordinary voicing rules (§\,\ref{sec:morphophon}).

\subsection{Quantitative classes and local grades}\label{sec:verbal-quantity}

Verbal quantitative class determines where an aspectual contrast receives
its quantity. The grades of §\,\ref{sec:morphophon} distinguish lexical
G0, strong G1, raised G2, and raised-strong G3. A grade selects one
licensed nucleus in the local stem, rather than lengthening every vowel
in a written word. The target ordinarily is the rightmost eligible
root nucleus, excluding fleeting vowels and material belonging to a
prefix or suffix; a lexical target specification takes precedence.
In \ird{zemac-} \trsl{run}, the target is its second vowel, /a/:
G1 gives \ird{zemác-}, preserving the initial /ɛ/ and final /ts/.
Silent \orth{i} in \ird{piašt-} contributes no nucleus; its strong
grade is \ird{piášt-}.

Class~A expresses the strong progressive and perfective in its root
grade, while the continuous, contemplative, and retrospective select
G0. Class~B instead uses strong aspectual exponents with an otherwise
short root. Class~C preserves an independently long lexical root
and selects short aspectual endings. The principal parts in
Table~\ref{tab:verbal-quantity-classes} illustrate these specifications;
the perfectives of \ird{zemac-} and \ird{knář-} are intransitive,
whereas ordinary transitive \ird{piašt-} includes patient orientation
\ird{-n-}.

\begin{table}[htbp]
\footnotesize
\caption{Principal parts of the three verbal quantitative classes.}
\label{tab:verbal-quantity-classes}
\begin{tabularx}{\textwidth}{L L L L L L}
\toprule
{\sc class/stem} & {\sc continuous} & {\sc progressive} & {\sc contemplative} & {\sc perfective} & {\sc retrospective} \\
\midrule
A \ird{zemac-} & \ird{zemaca} & \ird{zemácime} & \ird{zemacáž} & \ird{zemácek} & \ird{zemacáv} \\
B \ird{piašt-} & \ird{piašta} & \ird{piaštíme} & \ird{piaštáž} & \ird{piaštník} & \ird{piaštnáv} \\
C \ird{knář-} & \ird{knářa} & \ird{knářime} & \ird{knářaž} & \ird{knářek} & \ird{knářav} \\
\bottomrule
\end{tabularx}
\end{table}

Class~C \ird{strád-} likewise retains its lexical long vowel in
\ird{stráda}, \ird{strádime}, and \ird{strádaž}.

Copular, existential, and other functional predicates retain their listed
principal parts; surface shortness alone does not assign them to Class~A
(Chapter~\ref{ch:simple}).

The Class~B perfective has long \ird{-ék}, with bound \ird{-ík-},
where a bare perfective is independently licensed, and long \ird{-ník}
with ordinary patient orientation. Its progressive has \ird{-íme}
and its retrospective \ird{-áv}, or \ird{-náv} after orientation.
Class~A instead combines root G1 with short perfective
\ird{-ek}/\ird{-nik} and short progressive \ird{-ime}; its
retrospective has short-root G0 and long \ird{-áv}/\ird{-náv}.
The stative verbaliser's perfective \ird{-ik} likewise follows the
stem's quantitative class: \ird{mrócnik} has root G1 and short /ɪ/.

Valency derivation establishes a new local quantity domain before
Asp.Align attaches. Replacement antipassive Voice selects a strong
stem such as \ird{piášč-}; compact causative \ird{-r-} selects G1
on an eligible local stem, as in \ird{piáštr-}; and applicative
\ird{-éb-} retains its own long vowel. Such stems select short
adjacent aspectual exponents: \ird{piášček}, \ird{piáščav},
\ird{piáštrime}, \ird{piaštébaž}. A weak aspect does not undo
strength already supplied by the inner head. Additive antipassives,
including \ird{dalš-}, \ird{kělš-}, and \ird{lovž-}, preserve their
listed short root and short \ird{-ek}/\ird{-av}; Aspect does not
reach through that Voice head to create a root-strong grade.

The short \ird{-až} is a conditioned contemplative allomorph after
Class~C or a locally quantity-bearing valency stem. Short
\ird{-av}, \ird{-ime}, and \ird{-ek}/\ird{-ik} likewise belong
to their licensed local combinations. These choices do not impose
one long vowel per morphological word. Directive Mood, participant
nominalisers, relatives, and outer anchoring preserve their independently
specified quantities unless their construction explicitly selects
another allomorph. Thus \ird{piáštrím} may contain both the
causative-selected root quantity and the long imperative vowel.

Height alternation remains independent of quantity. Raising under G2
preserves the selected vowel's quantity, while an affirmative irrealis
with an exposed, eligible root can select raised-strong G3. Mid
\ird{e/ě} then gives \ird{í}, and \ird{o} gives \ird{ú}, where
that raised relation is licensed; low-vowel \ird{piašt-} selects
strong \ird{piášč-} without an invented raising of /a/.
The consonantal grade is interpreted separately at its own boundary.
Ordinary indicative Aspect may preselect its class-assigned root grade
before assembly with plain patient orientation \ird{-n-}, as in
Class~A \ird{vádnik}. This construction-specific selection does not
make orientation transparent to every later head. Mood instead acts on
the assembled local domain: intervening valency or orientation prevents
it from reaching back to assign a new root grade
(§\,\ref{sec:vm-moodp}).


\section{Aspect and alignment}
\label{sec:aspect}

Iridian grammaticalises \term{aspect} rather than \term{tense} on the finite
verb. There is no dedicated past-, present-, or future-tense suffix; temporal
location is inferred from discourse context, temporal adverbials, and the
lexical semantics of the predicate. Aspect and alignment are fused into a single
obligatory suffix, the \term{aspect-alignment morpheme} (Asp.Align), which every
finite verb must bear: selecting an aspectual category simultaneously determines
how the event is construed and the alignment domain within which
valency and orientation determine the privileged participant.

The five \synhead{Asp.Align} categories fall into two domains. The
\term{perfective domain} comprises the perfective and retrospective, which
present an eventuality as bounded or anterior and license the ordinary
ergative-absolutive frame: patient direct, agent transitive. The
\term{imperfective domain} comprises the continuous, progressive, and
contemplative, which present an eventuality internally or prospectively and
license the ordinary nominative-accusative frame: agent direct, patient
indirect. Applicatives, antipassives, and patient-selecting Voice modify these
defaults (§\,\ref{sec:vm-asp-case}). The case split belongs to Iridian's
grammar; temporal viewpoint alone does not determine the selected participant
or confer aboutness topichood.

The subsections that follow first lay out the shared temporal-semantic machinery
(including Reichenbachian framing and the formal Asp.Align analysis), then treat
the perfective domain (§\,\ref{sec:aspect-pf}) and the imperfective domain
(§\,\ref{sec:aspect-ipf}) in turn.

Three temporal notions must be kept distinct. Tense locates a situation in time
relative to a reference point, most commonly the moment of speaking. Grammatical
aspect does not locate the situation in this way, but construes its internal
temporal contour: whether it is viewed as bounded, ongoing, prior from a later
vantage point, or not yet realised. \term{Lexical aspect} or \term{Aktionsart},
by contrast, belongs to the predicate itself. A predicate may denote a state,
activity, accomplishment, or achievement independently of the inflection it
bears. Iridian aspectual morphology therefore does not determine the inherent
temporal type of the eventuality; rather, it imposes a viewpoint upon an
eventuality whose lexical contour is supplied by the verb phrase.

Lexical preverbs make this distinction overt: \ird{mněktař-} \trsl{read
through} specifies traversal of a text, while its continuous and progressive
forms describe that traversal without asserting arrival at its endpoint.
The same stem admits all five Asp.Align categories, with the ordinary case
frames and Voice requirements (§\,\ref{sec:preverb-aspect}). Likewise,
\ird{ranóčkime} \trsl{is running into an enclosure} describes an ingress
in progress without asserting that entry has been completed, while
\ird{rustójime} \trsl{is retracing a route} locates the returning motion
within its unfolding course. Direction, configuration, and restitution belong
to the lexical description; continuous, progressive, contemplative, perfective,
and retrospective morphology remain available wherever the resulting meaning
licenses them.

A useful descriptive frame is Reichenbach's distinction among speech time
(\textit{S}), reference time (\textit{R}), and event time (\textit{E}):
\textit{S} is the moment of utterance, \textit{E} the time at which the
situation itself takes place, and \textit{R} the temporal vantage point from
which that situation is presented. The usefulness of \textit{R} is that speakers
do not always describe an event directly from the moment of speaking, but may
instead locate it from some later or earlier point already established in the
discourse; English \textit{had left}, for example, places the leaving
(\textit{E}) before a reference point (\textit{R}) that is itself prior to the
speech moment (\textit{S}). In a tense-prominent system, finite morphology often
fixes the relation among all three. In Iridian, by contrast, the finite suffix
primarily construes the relation between \textit{E} and \textit{R}, while the
relation of \textit{R} to \textit{S} is often left to context, temporal
adverbials, and discourse sequencing. This is why the same aspectual category
may receive past-, present-, or future-oriented readings in different
environments without ceasing to be aspectual rather than temporal.

The five categories may be sketched in these terms. The \term{perfective}
presents an eventuality as bounded and viewed as a whole, with \textit{R}
aligned to the event in its entirety. The \term{retrospective} places \textit{E}
before \textit{R}, yielding a later vantage point from which present relevance,
result, prior experience, or other discourse effects of anteriority may be
construed. The \term{continuous} and \term{progressive} both place \textit{R}
within \textit{E}, but they differ in the kind of internal view imposed: the
continuous presents the situation as simply obtaining at reference time, without
foregrounding change, onset, or unfolding contour, whereas the progressive more
narrowly profiles a dynamic event as actively unfolding at the reference point.
The \term{contemplative} places \textit{E} after \textit{R}, presenting the
event as projected, intended, expected, or impending rather than as obtaining or
complete.

These five categories divide naturally into two families. The
\term{perfective domain}---the perfective and retrospective---consists of those
categories that locate the event from an external or posterior vantage and
present it as bounded or anteriorly packaged: in Reichenbachian terms, the
event time is enclosed within the reference span, or precedes it. The
\term{imperfective domain}---the continuous, progressive, and
contemplative---consists of those categories that do not present the event as a
bounded whole, but instead either locate reference time within the ongoing event
or situate reference time prior to an unrealized projection of the event. The
Reichenbachian figure (Figure~\ref{fig:iridian-aspect-reichenbach}) makes this
grouping visually explicit: the upper two rows belong to the perfective domain,
in which the event is viewed as bounded or anterior from outside; the lower
three rows belong to the imperfective domain, in which R lies within E or
precedes an unrealized E.

Because viewpoint aspect interacts with lexical aspect, the same inflection does
not yield identical interpretive effects with every predicate class. With states
and activities, the continuous naturally yields persistent or ongoing readings;
with accomplishments it typically selects an internal phase rather than a
completed endpoint; with achievements it may invite iterative, preparatory, or
inceptive coercions. Conversely, the perfective and retrospective readily
package telic predicates as bounded wholes or prior results, but with atelic
predicates they instead delimit a stretch, occasion, or state viewed from an
external vantage. The exact semantic and distributional consequences of these
interactions are discussed in the domain subsections below.

Typologically, this pattern is an instance of aspect-conditioned split
ergativity. Iridian does not maintain a single alignment pattern across all
finite clauses; rather, the aspectual category selected on the verb determines
whether the clause is built on an ergative-absolutive or a nominative-accusative
frame. Aspect in Iridian therefore has consequences at two levels at once: it
construes the temporal profile of the eventuality and it determines the
morphosyntactic alignment through which the clause is organised.

The aspectual domains condition case without encoding a discourse topic.
A perfective transitive may continue reference to its agent, although the
patient takes direct case, while an imperfective clause may introduce an event
without an established topic. Valency and orientation morphology determine
which participant receives direct-case privileges within the aspectual domain;
linked topicalisation independently determines aboutness
(§\,\ref{sec:vm-asp-case}; Chapter~\ref{ch:infstr}).

\begin{figure}[t]
\centering
\begin{tikzpicture}[
  x=1.05cm,y=1cm,
  evt/.style={line width=1.1pt},
  refpt/.style={circle,fill,inner sep=1.3pt},
  speech/.style={diamond,draw,inner sep=1.2pt},
  rowlab/.style={anchor=east,font=\small},
  every node/.style={font=\small}
]

% All rows hold S at x = 4.5 so the deictic baseline is visually constant.

% perfective (y = 0)
\node[rowlab] at (0,0) {perfective};
\draw[->] (0.5,0) -- (8.5,0);
\draw[evt] (3.0,0) -- (6.0,0);
\draw (3.0,0.15) -- (3.0,-0.15);
\draw (6.0,0.15) -- (6.0,-0.15);
\node[refpt] at (3.0,0) {};
\node[refpt] at (6.0,0) {};
% R bracket above the timeline, E label between bracket and line
\node at (4.5,0.30) {$E$};
\draw[decorate,decoration={brace,amplitude=3pt}]
  (3.0,0.55) -- (6.0,0.55)
  node[midway,above=4pt] {$R$};
% S below, matching spacing of other rows
\draw[dashed] (4.5,-0.08) -- (4.5,-0.42);
\node[speech] (s1) at (4.5,-0.55) {};
\node[below=2pt of s1] {$S$};

% retrospective (y = -2.0)
\node[rowlab] at (0,-2.0) {retrospective};
\draw[->] (0.5,-2.0) -- (8.5,-2.0);
\draw[evt] (1.8,-2.0) -- (3.4,-2.0);
\draw (1.8,-1.85) -- (1.8,-2.15);
\draw (3.4,-1.85) -- (3.4,-2.15);
\node at (2.6,-1.65) {$E$};
\node[refpt,label=above:{$R$}] at (4.5,-2.0) {};
\node[speech] (s2) at (4.5,-2.55) {};
\node[below=2pt of s2] {$S$};
\draw[dashed] (4.5,-2.08) -- (4.5,-2.42);

% continuous (y = -4.0)
\node[rowlab] at (0,-4.0) {continuous};
\draw[->] (0.5,-4.0) -- (8.5,-4.0);
\draw[dotted] (1.6,-4.0) -- (2.2,-4.0);
\draw[evt] (2.2,-4.0) -- (6.8,-4.0);
\draw[dotted] (6.8,-4.0) -- (7.4,-4.0);
\node at (3.0,-3.65) {$E$};
\node[refpt,label=above:{$R$}] at (4.5,-4.0) {};
\node[speech] (s3) at (4.5,-4.55) {};
\node[below=2pt of s3] {$S$};
\draw[dashed] (4.5,-4.08) -- (4.5,-4.42);

% progressive (y = -6.2)
\node[rowlab] at (0,-6.2) {progressive};
\draw[->] (0.5,-6.2) -- (8.5,-6.2);
\draw[evt] (2.4,-6.2) -- (7.0,-6.2);
\draw (2.4,-6.05) -- (2.4,-6.35);
\draw (7.0,-6.05) -- (7.0,-6.35);
\node at (6.5,-5.85) {$E$};
\draw[->] (2.8,-5.55) -- (4.2,-5.55)
  node[midway,above,font=\scriptsize] {dynamic phase};
\node[refpt,label=above:{$R$}] at (4.5,-6.2) {};
\node[speech] (s4) at (4.5,-6.75) {};
\node[below=2pt of s4] {$S$};
\draw[dashed] (4.5,-6.28) -- (4.5,-6.62);

% contemplative (y = -8.2)
\node[rowlab] at (0,-8.2) {contemplative};
\draw[->] (0.5,-8.2) -- (8.5,-8.2);
\node[refpt,label=above:{$R$}] at (4.5,-8.2) {};
\draw[evt] (5.8,-8.2) -- (7.5,-8.2);
\draw (5.8,-8.05) -- (5.8,-8.35);
\draw (7.5,-8.05) -- (7.5,-8.35);
\node at (6.65,-7.85) {$E$};
\node[speech] (s5) at (4.5,-8.75) {};
\node[below=2pt of s5] {$S$};
\draw[dashed] (4.5,-8.28) -- (4.5,-8.62);

\end{tikzpicture}
\caption{Reichenbachian schematic for Iridian aspect. $E$ = event time, $R$ =
reference time, $S$ = speech time. In the perfective, $R$ spans the bounded
event (marked by dots at the endpoints and a bracket); elsewhere, $R$ is a
point-like vantage. The upper two rows (perfective, retrospective) belong to the
\term{perfective domain}: in both, the event is construed as bounded or
anterior, viewed from outside or from a posterior vantage. The lower three rows
(continuous, progressive, contemplative) belong to the \term{imperfective
domain}: in all three, $R$ either lies within $E$ or precedes an unrealized $E$,
and no external terminus is foregrounded. Dashed connectors mark the unmarked
deictic baseline in which $R$ coincides with $S$. $S$~is held at a constant
horizontal position across all rows; in discourse, adverbials or sequencing may
shift $R$ independently of $S$.}
\label{fig:iridian-aspect-reichenbach}
\end{figure}

\subsection{Formal analysis: the Asp.Align morpheme}\label{sec:vm-asp}

The obligatory aspect-alignment suffix encodes viewpoint and the alignment
domain within which the verbal construction licenses case. The five indicative
exponents are given in Table~\ref{tab:aspect}.

\begin{table}[ht]
	\footnotesize
	\caption{Aspect-alignment exponents by domain and quantitative class.}\label{tab:aspect}
	\begin{tabularx}{0.98\textwidth}{L L L L}
		\toprule
		{\sc domain} & {\sc aspect} & {\sc exponent} & {\sc quantity and orientation} \\
		\midrule
		Perfective & Perfective & \ird{-ek}/\ird{-ék} & bare forms require an intransitive reading; ordinary transitive \ird{-nik}/\ird{-ník}; class-selected short or long bound \ird{-ik-}/\ird{-ík-} \\
		& Retrospective & \ird{-áv}/\ird{-av} & ordinary A/B long \ird{-áv}, or \ird{-náv} with orientation; C and listed strong-valency stems short \ird{-av}/\ird{-nav} \\
		\midrule
		Imperfective & Continuous & \ird{-a} & local G0 in ordinary A/B; previously selected lexical or valency strength retained \\
		& Progressive & \ird{-ime}/\ird{-íme} & A root G1 with short ending; B suffix-strong; C and strong-valency stem with short ending \\
		& Contemplative & \ird{-áž}/\ird{-až} & A/B ordinary G0 with long ending; C and strong-valency stem with short ending \\
		\bottomrule
	\end{tabularx}
\end{table}

Antipassive and applicative forms preserve the orientation and quantity
selected by their own constructions. Replacement antipassives such as
\ird{piášč-} combine with short \ird{-ek}/\ird{-av}, while the
applicative \ird{-éb-} selects short perfective \ird{-ik}.
The distribution in Table~\ref{tab:verbal-quantity-classes} supplies the
quantity distinctions without changing the five viewpoint categories.

\subsubsection{Split-ergative alignment}
\label{sec:vm-asp-case}

In finite indicative clauses, Asp.Align encodes viewpoint aspect and the
alignment domain. Together with valency and orientation morphology, it
determines the clause's \term{absolutive pivot}, the argument that bears
direct case. Its frequent interpretation as an
aboutness topic belongs to the independent information-structural system
(Chapter~\ref{ch:infstr}). Aspect and alignment are jointly realised in the obligatory
suffix, with case valuation active in the finite domain. Participant
nominalisation preserves the aspectual and participant-selection features
without finite case valuation (§\,\ref{sec:nominalised}).

The ordinary active frames distinguish two domains. In the ERG-aspect domain
(perfective and retrospective), the patient takes direct case and the agent
transitive case; ordinary transitives require \ird{-n-}, whereas an
inherently intransitive predicate or the intransitive reading of an
ambitransitive takes the bare aspect suffix. In the NOM-aspect domain
(continuous, progressive and contemplative), the agent takes direct case and
the patient indirect case. The patient remains a core argument in either frame.

Valency morphology modifies the argument inventory before orientation and case
are determined. A simple antipassive licenses only its actor as an independent
core participant, and that actor takes direct case. A simple applicative
selects the applied participant as direct; its actor takes indirect case in the
NOM domain and transitive case in the ERG domain, while any retained theme is
indirect. A higher causative establishes a new orientation domain, in which the
causer is direct in NOM aspects; in ERG aspects the selected participant is the
retained lexical theme, otherwise the applied participant, otherwise the
causee (§\,\ref{sec:vm-caus}). These are constructional extensions of the
ordinary split, rather than additional meanings of viewpoint aspect.

The alignment split is an established property of Iridian, whose aspectual
 viewpoint and case assignment are jointly realised in Asp.Align. Perfective
 and retrospective transitives with \ird{-n-} assign direct case to the patient
 and transitive case to the agent, while the imperfective active frame assigns
 direct case to the agent and indirect case to the patient.
 These orientations concern grammatical realisation: a bounded event does not
 semantically entail a patient topic, and an ongoing activity does not require
 an agent topic. The grammar therefore treats their discourse association as
 a default rather than deriving aboutness from event culmination.

Event construal and information structure consequently make distinct
 contributions: an event may be presented as completed while the utterance
 continues a discussion of its agent, just as an ongoing event may be
 introduced without establishing any participant as an aboutness topic.

The ordinary active frames illustrate this contrast.
In the NOM-aspect domain, the agent or sole argument bears direct case,
while the patient of an ordinary transitive appears in
the indirect. This is a nominative-accusative frame, with the direct case
marking the privileged participant and the indirect marking the core patient:

\pex
\a
\begingl
\gla Marek zemácime.//
\glb Marek.\Dir{} run-\Prog{}//
\glft \trsl{Marek is running.} \lingcontext{intransitive}//
\endgl
\a
\begingl
\gla Marek travě piaštíme.//
\glb Marek.\Dir{} bread.\Ind{} eat-\Prog{}//
\glft \trsl{Marek is eating bread.} \lingcontext{transitive; P = \Ind{}}//
\endgl
\a
\begingl
\gla Nužo mrocna.//
\glb city.\Dir{}.\Def{} be.beautiful-\Cont{}//
\glft \trsl{The city is beautiful.} \lingcontext{stative continuous}//
\endgl
\a
\begingl
\gla Nužo mrócnime.//
\glb city.\Dir{}.\Def{} be.beautiful-\Prog{}//
\glft \trsl{The city is becoming beautiful.} \lingcontext{dynamic progressive}//
\endgl
\xe

In an ordinary ERG-domain transitive, by contrast, the patient is
placed in the direct case, and the agent takes the transitive case. This is an
ergative-absolutive frame; what Iridian grammar calls the \emph{transitive case}
is, in typological terms, the ergative, and the \emph{direct case} serves as
the absolutive. An ordinary transitive in this domain requires the orientation head
\ird{-n-} to select its patient; when antipassive
\ird{-aš-} is present instead, the agent takes the absolutive and the patient is
removed from core argument structure (§\,\ref{sec:vm-antipassive}):

\pex
\a
\begingl
\gla Travó Marku piaštník.//
\glb bread.\Dir{}.\Def{} Marek.\Trs{} eat-\textsc{trs}-\Pf{}//
\glft \trsl{Marek ate the bread.} \lingcontext{ergative; P = \Dir{}}//
\endgl
\a
\begingl
\gla Marek piášček.//
\glb Marek.\Dir{} eat-\Antip{}-\Pf{}//
\glft \trsl{Marek ate.} \lingcontext{antipassive; A = \Dir{}}//
\endgl
\xe

The ergative transitive licenses a core patient in direct case, whether that
patient is overtly expressed or recovered from context. In a narrative about
that participant, the clause can describe what happened to it, as in the
bread example above; the case frame also remains available when the agent is
independently topicalised (§\,\ref{sec:tc-marked-topics}). Suffixal anchoring
identifies an accessible referent without being required by the construction.
Neither direct case nor transitive Voice entails that the patient undergoes
a change of state: the construction equally accommodates predicates of buying
and perception.

The antipassive presents the agent's activity without independently licensing
a patient, although the lexical event may still involve a conceptual theme.
This frame suits activity-oriented contexts in which no referential patient is
asserted or recovered. It excludes an indirect-case patient as well as a
direct-case one (§\,\ref{sec:vm-antipassive-desc}); an ordinary transitive
with a contextually recovered patient retains a different argument structure.

The ergative transitive and antipassive differ in the realisation of the
 patient and the selection of the absolutive pivot. The former licenses a
 core patient; the latter removes the patient from core argument structure.
 Their discourse uses follow from those valency differences, while aboutness
 remains independently available to suitable referents.

Iridian's alignment is \term{aspect-directed}: ordinary clauses distinguish
patient-direct ERG frames from agent-direct NOM frames. Applicative and
patient-selecting constructions establish their own orientation within those
domains, so the aspect value is one condition on the case frame rather than a
complete specification of its pivot.

The transitive case of an ordinary ERG-domain agent is analysed here as
dependent on a distinct participant licensed in direct case. Applicatives and
causatives extend this relation to their constructionally selected participant,
which need not be the lexical patient. This dependent-case analysis describes
the Iridian mapping; the surface contrast alone does not determine a universal
algorithm for case assignment.

Aspect selection establishes a case frame and temporal viewpoint, whose
 interaction can favour different discourse continuations. The pivot is the
 default aboutness candidate where the context supports ordinary topical
 predication, but a speaker can continue an agent topic in a perfective
 transitive without changing the event's aspect. The linked-topic dependency
 retains the agent's transitive case and the patient's direct case
 (§\,\ref{sec:tc-marked-topics}).

The suffix \ird{-n-} realises the higher orientation head, outside causative
morphology where present and taking \ird{-an-} after a rhotic causative.
It selects the patient of an ordinary transitive or
the highest eligible non-causer of a derived causative. In the ERG domain it
combines with the ordinary transitive frame; in the NOM domain its
patient-oriented use licenses a passive, with the agent restricted to optional
indirect expression (§\,\ref{sec:vm-passive}). Simple applicatives establish
their own orientation through \ird{-éb-}, without a second overt \ird{-n-}.

Antipassive morphology builds a reduced predicate before case is assigned.
Its occurrence is licensed by the completed ERG-aspect configuration; it does
not inspect a patient whose case has already been valued. A following
causative may introduce a new causer and an outer \ird{-an-}, although the
lexical patient remains suppressed (§\,\ref{sec:vm-antipassive}).

\subsection{The perfective domain}\label{sec:aspect-pf}

The perfective domain comprises those Asp.Align forms which construe an
eventuality as bounded and externally viewed. Unlike the imperfective domain,
which presents an event from within its internal course, the perfective domain
presents it as a temporally delimited whole. Iridian distinguishes two
categories within this domain: the \term{perfective} proper, which presents a
bounded event as a discrete occurrence, and the \term{retrospective}, which
presents a bounded prior event in relation to a present state, present
relevance, or present experience. Both categories interact regularly with
valency and alignment morphology, and both therefore belong to a single
structural and semantic domain rather than to unrelated aspectual classes.

Free perfective quantity follows the selected quantitative class.
Class~A \ird{kélek} \trsl{sang} and \ird{ščének} \trsl{arrived}
have root G1 and short \ird{-ek}; Class~B has long \ird{-ék} where
a bare perfective is licensed. Adjacent linker \ird{-e} and participant
nominaliser \ird{-ou} select the corresponding bound height grade,
short \ird{-ik-} or long \ird{-ík-}. A replacement antipassive
already supplies root strength: \ird{piášček}, linked
\ird{piáščice}, and participant \ird{piáščikou} therefore
retain short suffix vowels (§\,\ref{sec:pf-dissim}).

Ordinary patient orientation combines \ird{-n-} with the class-selected
high-vowel perfective: short \ird{-nik} in root-strong Class~A and
long \ird{-ník} in suffix-strong Class~B. The stative verbaliser
\ird{-n-} independently selects short \ird{-ik}, with root quantity
assigned by its stative class; its surface \ird{-nik} does not insert
transitive Voice into the stative predicate. The selected quantities
persist before further morphology. Retrospective likewise distinguishes
ordinary A/B \ird{-áv}/\ird{-náv} from C and strong-valency
\ird{-av}/\ird{-nav}; antipassive \ird{piáščav} has a strong
replacement stem and short ending.

The semantic opposition within the perfective domain is best stated as one of
\emph{event occurrence} versus \emph{post-event relevance}. The perfective
proper presents a bounded event as a single occurrence situated at a definite
point in narrative or discourse time, and is therefore the primary narrative
aspect. The retrospective likewise presupposes a bounded prior event, but does
not profile that event merely as an occurrence; rather, it presents it through
its present relevance, present consequences, or experiential force. In this
respect it corresponds in some uses to the English perfect. Its non-volitional
or accidental readings belong to the same general semantic profile: the event is
construed from the standpoint of its aftermath or present bearing, rather than
simply narrated as a bounded occurrence.

For this reason, speakers often choose the retrospective where reduced
volitionality is pragmatically inferable or interactionally relevant. By
presenting the event through its resulting situation rather than as a discrete
occurrence in the narrative line, the retrospective readily backgrounds agency
and favours readings such as `it happened', `it turned out that way', or `I find
myself in the state resulting from it'. This preference is pragmatic rather
than semantic: the retrospective does not itself encode non-volitionality, but
it is frequently selected in contexts where non-control, accident, or
unintended outcome is to be understood.

This contrast may be seen in the following pair:

\pex
\a
\begingl
\gla Praže de stožáv.//
\glb Prague.\Ind{} to go-\Ret{}//
\glft \trsl{(I) have been to Prague before.}//
\endgl
\a Compare:\par
\begingl
\gla Bych Praže de stóžek.//
\glb yesterday Prague.\Ind{} to go-\Pf{}//
\glft \trsl{Yesterday I went to Prague.} \lingcontext{specific past event; perfective}//
\endgl
\xe

Within the perfective domain, the ordinary transitive selects its patient
through \ird{-n-}, assigning direct case to that participant and transitive
case to the agent. An antipassive instead leaves only the actor independently
licensed, whereas a simple applicative selects its applied participant. An
inherently intransitive verb requires no overt orientation marker, since its
sole argument takes direct case (§\,\ref{sec:vm-asp-case}).

The retrospective can present a patient-related outcome in an ergative
transitive or an agent's prior experience in an antipassive. The latter has no
independent patient argument whose resulting state could supply the discourse
centre, although the activity's consequences may remain relevant. These are
contextual interpretations of the respective argument frames, as illustrated
by English \textit{the bread has been eaten} and \textit{I have eaten};
retrospective morphology itself does not require either a result-state or an
experiential reading.

Negation further confirms the semantic distinction between the two members of
the perfective domain. When negated, the perfective normally carries the
implicature that the event was expected, intended, or contextually projected,
but failed to occur. The retrospective, by contrast, is the neutral form for
simple past negation without this implication. The opposition is illustrated
below:

\pex
\a
\begingl
\gla Marek zapiášček.//
\glb Marek.\Dir{} \Neg{}-eat-\Antip{}-\Pf{}//
\glft \trsl{Marek didn't eat (but he should have).}//
\endgl
\a
\begingl
\gla Marek zapiáščav.//
\glb Marek.\Dir{} \Neg{}-eat-\Antip{}-\Ret{}//
\glft \trsl{Marek didn't eat (neutral past).}//
\endgl
\xe

\label{sec:ambitransitive}A small class of verbs shows particularly clearly how the perfective domain
interacts with argument structure. These are lexical ambitransitives: verbs
which possess both an intransitive reading, in which no patient is structurally
available, and a transitive reading, in which a patient is licensed. Neither
reading is derived from the other by ordinary valency operations; both are
lexical possibilities of the same stem. Two examples are \ird{lov-} \trsl{hunt,
forage [intr.]; catch, take [tr.]} and \ird{kěl-} \trsl{sing, perform [intr.];
sing a song, chant a text [tr.]}.

The diagnostic for the two readings is whether ordinary transitive orientation \ird{-n-}
is available in one reading but impossible in the other. In \ird{lov-}, the
intransitive perfective---formed without any Voice head, with the ERG Asp.Align
suffix attached directly to the stem---accesses only the intransitive reading:
the event is construed simply as hunting, with no patient structurally present
and no quarry included in the event frame as an affected participant. The
ergative and antipassive forms, by contrast, access the transitive reading. The
ergative perfective \ird{lóvnik} is possible only when something is caught or
taken; the antipassive construes the event as directed toward prey, but
suppresses that prey from the clause's core argument structure.

\pex
\a
\begingl
\gla Marek bych lóvek.//
\glb Marek.\Dir{} yesterday hunt-\Pf{}//
\glft \trsl{Marek went hunting yesterday.} \lingcontext{intransitive reading; no patient in event frame}//
\endgl
\a
\begingl
\gla Zvěrato Marku lóvnik.//
\glb game.\Dir{}.\Def{} Marek.\Trs{} hunt-\textsc{trs}-\Pf{}//
\glft \trsl{Marek caught the game.} \lingcontext{transitive reading; ergative}//
\endgl
\a
\begingl
\gla Marek bych lovžek.//
\glb Marek.\Dir{} yesterday hunt-\Antip{}-\Pf{}//
\glft \trsl{Marek was hunting yesterday.} \lingcontext{transitive reading; antipassive; quarry conceptually present but not licensed}//
\endgl
\xe

The two non-ergative forms are therefore not interchangeable. \ird{Lóvek}
belongs to the intransitive lexical reading: hunting is construed as an activity
with no patient present even in the event structure. \ird{Lovžek}, however,
belongs to the transitive lexical reading: the event is directed toward prey,
and that prey remains part of the conceptual frame even though the antipassive
removes it from core syntax. The distinction is thus not merely one of overt
argument expression, but of underlying event construal.

The stem \ird{kěl-} likewise distinguishes an ordinary intransitive perfective,
a transitive perfective, and an antipassive:

\pex
\a \phon{kélek}{ˈkʲeː.lɛk}{sang, performed} \lingcontext{intransitive; no song in the event frame}
\a \phon{kélnik}{ˈkʲeːw.nɪk}{sang it, sang a specific song} \lingcontext{transitive; ergative; patient licensed}
\a \phon{kělšek}{ˈkʲɛw.ʃɛk}{was singing} \lingcontext{transitive reading; antipassive; patient suppressed from core syntax}
\xe

In \ird{kélek}, ordinary \ird{-ek} follows the full stem without root-height
raising: G1 lengthens /ɛ/ to /eː/. Transitive \ird{kélnik} contains Voice \ird{-n-} and the
construction-specific high-vowel perfective; the difference between \ird{e} and
\ird{i} belongs to that exponent selection. Antipassive \ird{kělšek} takes
a short free perfective after its augmented stem. Class~A root G1
supplies long \ird{é} in \ird{kélek} and \ird{kélnik}, whereas
additive Voice in \ird{kělšek} retains short \ird{ě}. The forms
distinguish valency and event construal together with their licensed quantity
patterns.
Their initial consonant is lexical /kʲ/, whose velar closure and secondary
palatal articulation are retained in all three forms; it is distinct from
\ird{ť} [c] (Table~\ref{tab:kt-contrast}). The short-grade spelling \ird{kě} and long-grade \ird{ké} retain
the lexical /kʲ/ reading; quantity selection supplies no additional
consonant mutation.
The lateral remains [l] in the onset of \ird{kélek}, but surfaces as coda
[w] in \ird{kélnik} and \ird{kělšek}, according to the ordinary rule
(§\,\ref{sec:l-vocalisation}).

\subsection{The imperfective domain: internal and projected viewpoint}
\label{sec:aspect-ipf}

The imperfective domain comprises those Asp.Align categories that do not
construe the eventuality as a bounded whole viewed from outside. Where the
perfective domain packages an event as a closed unit and foregrounds its
culmination or its anteriorly packaged result, the imperfective domain takes
no such external vantage: it either locates the reference point within the
ongoing event, without presenting its boundaries as salient, or presents the
event as projected relative to reference time---not yet realised, but
anticipated, intended, or impending. These temporal properties characterise
the domain independently of its ordinary agent-direct case frame.

Three categories constitute the domain. The \term{continuous} is the unmarked
imperfective member; it presents a situation as simply obtaining at reference
time, without emphasis on change, onset, dynamic unfolding, or endpoint. The
\term{progressive} is the marked dynamic member; it profiles an event
specifically in its active unfolding phase. The \term{contemplative} is the
projective member; it presents the event as not yet realised at reference time
but projected from that point---intended, expected, impending, planned, or
otherwise prospectively construed. Together the three define an internal
opposition within the domain: static holding versus dynamic unfolding versus
prospective projection.

\paragraph{Alignment in the imperfective domain.}
The three imperfective categories license the ordinary NOM-domain frame,
with an agent or sole participant in direct case and a retained patient in
indirect case. A simple applicative instead selects the applied participant,
whereas overt \ird{-n-} selects the patient and licenses a passive. These
contrasts show why the temporal profile shared by the three aspects cannot by
itself determine which participant is privileged (§\,\ref{sec:vm-asp}).

\paragraph{The continuous (\ird{-a}).}
The continuous is the unmarked member of the imperfective domain. It presents
a situation as obtaining at reference time without foregrounding change, onset,
active unfolding, or prospective projection. It is therefore the natural form
for stable states, contingent present conditions, broad durative situations, and
habitual or generic activity where what matters is that the situation holds
rather than how it is developing:

\pex
\a
\begingl
\gla Praže to mulša.//
\glb Prague.\Ind{} \textsc{loc} live-\Cont{}//
\glft \trsl{(I) live in Prague.} \lingcontext{stable durative state}//
\endgl
\a
\begingl
\gla Nužo mrocna.//
\glb city.\Dir{}.\Def{} be.beautiful-\Cont{}//
\glft \trsl{The city is beautiful.} \lingcontext{property holding at the reference time}//
\endgl
\xe

Because the continuous does not foreground development or change, it is
compatible with predicates across the entire stative-dynamic spectrum, from
permanent properties to ongoing activities. Where the speaker wishes simply to
assert that a situation holds---without committing to whether it is in active
motion or stable---the continuous is the default choice. Its unmarked status
within the imperfective domain is reflected in the morphology: \ird{-a} is the
shortest and phonologically least specified of the three imperfective suffixes.

\paragraph{The progressive (\ird{-ime}).}
The progressive is the marked dynamic member of the imperfective domain.
Where the continuous merely asserts that a situation holds, the progressive
profiles an event specifically in its active unfolding phase---the period during
which it is dynamically occurring, not yet complete, and currently in motion.
It is accordingly most natural with activities and developing processes, where
the sense of an event in progress is pragmatically salient:

\pex
\a
\begingl
\gla Marek zemácime.//
\glb Marek.\Dir{} run-\Prog{}//
\glft \trsl{Marek is running.} \lingcontext{dynamic activity in progress}//
\endgl
\a
\begingl
\gla Marek travě piaštíme.//
\glb Marek.\Dir{} bread.\Ind{} eat-\Prog{}//
\glft \trsl{Marek is eating bread.} \lingcontext{ongoing activity; patient in indirect}//
\endgl
\xe

The progressive is more marked than the continuous with stative predicates.
Where the continuous with a stative stem yields simple present-state
predication, the progressive imposes a developmental or ingressive
interpretation---the state is not merely current but is in the process of coming
to obtain. The minimal pair already shown in the alignment section illustrates
this:

\pex
\a
\begingl
\gla Nužo mrocna.//
\glb city.\Dir{}.\Def{} be.beautiful-\Cont{}//
\glft \trsl{The city is beautiful.} \lingcontext{continuous: state holds at $R$}//
\endgl
\a
\begingl
\gla Nužo mrócnime.//
\glb city.\Dir{}.\Def{} be.beautiful-\Prog{}//
\glft \trsl{The city is becoming beautiful.} \lingcontext{progressive: dynamic unfolding at $R$}//
\endgl
\xe

This contrast---between the continuous as the form of a holding situation and
the progressive as the form of an actively unfolding process---is the most
consequential internal opposition in the imperfective domain. The behaviour of
stative predicates provides especially sharp evidence for it and is examined in
detail in §\,\ref{sec:stative-stability}.

\paragraph{The contemplative (\ird{-áž}).}
Default \ird{-áž} represents /aːʒ/ in ordinary short-root
Classes~A and~B. Class~C and locally quantity-bearing valency stems
select short \ird{-až}, as in \ird{knářaž} and \ird{piáštraž}.
Within each selected allomorph its quantity is retained before further
suffixes, while its final consonant has [ʃ] in an isolated
word by regular final devoicing. Before a following vowel, /ʒ/ retains
voicing: linked \ird{mrocnáže}, evidential \ird{piaštážal}, and
participant nominal \ird{piaštážou} all contain [ʒ]. The linker
preserves this postalveolar consonant rather than selecting a different
stem grade.

The contemplative is the projective member of the imperfective domain. It
presents the eventuality as not yet realised at reference time but projected
from that point: the event is construed as intended, expected, impending,
planned, or otherwise prospectively oriented. The contemplative therefore does
not encode future tense---like all Iridian Asp.Align categories it is
indifferent to the relation between reference time and speech time---but rather
a prospective viewpoint in which reference time is interior to the projection
and the event lies ahead of it:

\pex
\a
\begingl
\gla Praže de stojáž.//
\glb Prah-ě de stoj-áž.//
\glc Prague.\Ind{} to go-\Ctpl{}//
\glft \trsl{(I) am about to go to Prague. / (I) intend to go to Prague.}//
\endgl
\a
\begingl
\gla Nužo mrocnáž.//
\glb city.\Dir{}.\Def{} be.beautiful-\Ctpl{}//
\glft \trsl{The city will be beautiful. / The city is expected to be beautiful.}//
\endgl
\xe

The projective semantics account for the range of English translations the
contemplative receives---\textit{about to}, \textit{intending to},
\textit{planning to}, or a simple future---depending on predicate type and
discourse context. With achievement predicates, which denote instantaneous
transitions, the contemplative is especially natural: instantaneous transitions
lend themselves readily to imminence and anticipation construals, where the
event is imminent or impending rather than in active progress.

The contemplative belongs to the imperfective domain because reference time
precedes the projected eventuality, whose completion is unasserted. An ordinary
active contemplative selects its agent, or its sole participant with a stative
or intransitive base; this NOM-domain orientation remains subject to the same
valency and Voice alternations as the continuous and progressive.

The three imperfective categories interact systematically with predicate class,
showing that they are distinct viewpoint operators rather than translation
variants of a single value. With {states}, the continuous gives ordinary
present-state predication; the progressive is marked and typically forces a
developmental or ingressive reading (\trsl{getting better}, \trsl{becoming
\ldots}); the contemplative gives an anticipated or expected-state reading
(\trsl{will be \ldots}, \trsl{is expected to be \ldots}). With {activities}, the
continuous gives broad ongoing or habitual interpretation; the progressive gives
a vivid, scene-internal reading in which the activity is actively in course; the
contemplative gives an intended or pending-activity reading. With
{accomplishments}, the continuous selects the internal phase without projecting
completion; the progressive highlights active advancement toward culmination;
the contemplative presents the event as pending or aimed at, with the endpoint
anticipated but not asserted. With {achievements}, which denote instantaneous
transitions, the continuous and progressive typically require coercion into
preparatory, iterative, incipient, or extended interpretations, while the
contemplative is often especially natural because instantaneous transitions
readily support imminence readings.

Imperfective clauses permit descriptions of an agent's ongoing activity,
condition, or intentions, although their use does not require an agent topic.
The three categories support distinct discourse functions. The continuous is the natural form for description and
framing---durable predications, stable conditions, and ongoing situations that
provide the setting or context for dynamic foreground events. It is the unmarked
imperfective of persistence and background. The progressive is suited to
scene-internal narration, vivid observation, and dynamic immediacy---contexts
where an event is presented as actively unfolding before the narrator's or
speaker's attention. It is characteristically a foregrounded imperfective: not
the background frame but the observed process. The contemplative serves
planning, expectation, prefiguration, and other forward-looking discourse
functions---expressing what an agent intends, what is expected to follow, or
what is imminent in the developing narrative.

Imperfective categories are especially compatible with phasal modification
because they present the event internally or prospectively rather than as a
closed whole, leaving the event's phasal boundaries available for further
specification. The particle \ird{so} in its \trsl{still} reading combines
naturally with the continuous to assert the continuation of a situation.

The contemplative, presenting an event as prospective, is correspondingly
natural with \ird{so} in its \trsl{not yet} reading, describing a situation that
remains unrealized but is projected to obtain. The full phasal particle system
is described in Chapter~\ref{ch:function-words}; what matters here is that the
imperfective domain's open, non-bounded temporal profile is what licenses phasal
modification of this kind.

\paragraph{Evidence from stative predicates.}
Stative predicates furnish especially clear evidence for the contrast between
continuous and progressive, because with statives the interpretive difference
between the two is categorical. With a stable property predicate, the continuous
yields simple present-state predication: \ird{Nužo mrocna} \trsl{the city is
beautiful}. The progressive with the same predicate forces a developmental
reading: \ird{Nužo mrócnime} \trsl{the city is becoming beautiful}. The two
are not interchangeable. With temporary bodily or experiential conditions the
contrast is equally clear: the continuous describes the present condition as
simply obtaining (\trsl{is tired}), while the progressive more readily yields an
incipient or changing-state reading (\trsl{is getting tired}, \trsl{is starting
to feel sick}). The stative contrast instantiates the same imperfective opposition, with
lexically conditioned developmental and ingressive interpretations
(§\,\ref{sec:stative-stability}).


\subsection{Asp.Align under non-indicative Mood: impoverishment and reduced opposition}
\label{sec:vm-mood-asp}

Every finite Iridian verb contains Asp.Align, whose exponent can be zero.
Non-indicative Mood selects a reduced aspectual configuration in syntax,
retaining an IPFV/PFV opposition together with the construction's orientation
requirements.

In the indicative, Asp.Align realises the full five-way viewpoint inventory
(progressive, continuous, contemplative, perfective, retrospective) with the
associated split-ergative alignment. The non-indicative moods---imperative,
prohibitive, hortative, and irrealis---select the coarser
\term{reduced Asp.Align} opposition shown in
Table~\ref{tab:reduced-asp}.

\begin{table}
  \footnotesize
  \caption{Reduced Asp.Align exponents in non-indicative clauses.}
  \label{tab:reduced-asp}
  \begin{tabularx}{0.9\textwidth}{l l L}
    \toprule
    \textsc{domain} & \textsc{exponent} & \textsc{semantic range} \\
    \midrule
    Imperfective (IPFV) & $\varnothing$ & unbounded, ongoing, habitual, generic;
                                        processual commands and optatives \\
    Perfective (PFV)    & \ird{-ov-}/\ird{-óv-} & bounded, single whole event, completive;
                                        ingressively directive \\
    \bottomrule
  \end{tabularx}
\end{table}

Reduced Asp.Align retains the two aspectual domains, while the clause's
orientation also depends on Mood. In the irrealis, ordinary IPFV clauses retain
agent-direct and patient-indirect case, and PFV clauses retain patient-direct
and agent-transitive case. Directive Mood imposes theme orientation in both
domains (§\,\ref{sec:mood-directive-case}). The reduced temporal opposition
therefore does not require every non-indicative construction to retain the
ordinary indicative case frame.

\paragraph{Morphological composition.}
The surface template for a non-indicative finite verb is therefore:

\begin{center}
\Stem{} + (\Appl{}) + (\Caus{}) + (\textsc{orientation}) + \textsc{reduced Asp.Align} + \textsc{mood}
\end{center}

\noindent where the reduced Asp.Align is silent ($\varnothing$) in the IPFV
domain and has \ird{-ov-} in the PFV domain, with locally strong
\ird{-óv-} selected before affirmative irrealis. Directive Mood retains
\ird{-ov-}. Representative paradigms
for the stem \ird{piašt-} `eat':

\begin{center}
\footnotesize
\begin{tabular}{l l l l}
\toprule
\textsc{domain} & \textsc{imperative} & \textsc{prohibitive} & \textsc{hortative} \\
\midrule
IPFV ($\varnothing$) & \ird{piaštím}   & \ird{piašté}    & \ird{piaštká}   \\
PFV (\ird{-ov-})   & \ird{piaštovím} & \ird{piaštové}  & \ird{piaštovká} \\
\bottomrule
\end{tabular}
\end{center}

The PFV exponent \ird{-ov-} preserves the thematic consonant of the lexical
root. When irrealis \ird{-ěv} follows, its yod-grade instruction targets
the final /v/ of the aspect exponent, whose palatal realisation is [ʋʲ],
while the intervening exponent prevents mutation of the root-final cluster.
The resulting form is \ird{piaštóvěv}, rather than \ird{*piáščóvěv};
the absence of root mutation is compatible with palatality at the suffixal
boundary (§\,\ref{sec:suffix-classes}).

The IPFV imperative encodes an ongoing,
habitual, or processual directive (the event is construed as
unbounded), while the PFV imperative encodes a bounded, completive
directive (the event is construed as a single whole). In the irrealis,
the IPFV domain covers generic non-actual scenarios and standing
hypotheticals; the PFV domain covers single bounded counterfactuals.

In the Distributed Morphology implementation, Vocabulary Insertion maps the
coarse IPFV/PFV distinction to $\varnothing$ or \ird{-ov-}. A PF
impoverishment rule can exclude the finer indicative exponents in this
environment, although it does not supply the syntactic reduction, determine
case, or remove semantic distinctions after interpretation. The
exponence restriction is:

\medskip
\noindent\begin{tabular}{@{}l@{\quad}l}
$[\textsc{Mood}: \text{non-indicative}]$ & $\Rightarrow$ realise only the coarse
aspectual contrast at Asp.Align; \\
& \hspace*{1.5em}retain only $[\pm\textsc{pfv-domain}]$.
\end{tabular}
\medskip

Mood is structurally higher than Asp, with intervening Neg where present
(Table~\ref{tab:vm-full-structure}). Its syntactic selection establishes the
reduced configuration before case valuation, while its exponent follows the
zero or overt Asp.Align exponent. Morphological neutralisation therefore
presupposes the licensed configuration rather than deriving its case frame.

Directive Mood imposes an orientation requirement on the reduced aspectual
domain: a retained theme takes direct case, regardless of the IPFV/PFV value.
Where no theme is retained, selection falls on an applied participant, then a
causee, or the sole argument of an underived intransitive. The agent or causer
of a clause with a distinct direct participant takes transitive case, whether
expressed or contextually understood. This requirement is specific to the
directive construction; omission of a subject does not transfer its case to
another argument (§\,\ref{sec:mood-directive-case}).

\section{Valency operations}

Three derivational operations may be applied to the verb stem before the
aspect-alignment suffix: the antipassive, the applicative, and the causative.
Where more than one is present, the order is fixed: antipassive or applicative
precedes the causative.

\subsection{The antipassive (\ird{-aš-})}
\label{sec:vm-antipassive-desc}

The antipassive \ird{-aš-} forms an agent-oriented activity predicate from a
transitive base without licensing an independent patient DP. In the perfective
and retrospective, its agent receives direct case; the patient remains part of
the conceptual interpretation of the activity but cannot be expressed within
the finite clause. Iridian restricts this construction to the ERG-aspect domain.
A transitive base in that domain consequently requires an overt derivation:
\ird{-n-} retains the patient, \ird{-aš-} suppresses its independent
realisation, and \ird{-éb-} selects an applied participant. Bare perfectives
belong to independently intransitive readings (§\,\ref{sec:ambitransitive}).

The traditional notation \ird{-aš-} identifies this Voice construction;
it does not require an invariant string \ird{aš} to be concatenated with
every stem. Voice is realised through the replacement or augmentation specified
in Table~\ref{tab:verbclasses}. Thus \ird{piašt-} selects \ird{piášč-},
whereas the additive classes retain their thematic consonant before the
sibilant. These are morphological exponents of antipassive Voice, independently
of PAL and of the velar selector used by other constructions. Any accompanying
root-vowel quantity alternation, such as \ird{vad-}/\ird{váž-}, is specified
with its vocalic grade rather than inferred from the consonantal replacement.
The stem \ird{piášč-} combines with perfective \ird{-ek}:

\pex
\a
\begingl
\gla Marek piášček.//
\glb Marek.\Dir{} eat-\Antip{}-\Pf{}//
\glft \trsl{Marek ate.}//
\endgl
\a
\begingl
\gla Marek houlam se piášček.//
\glb Marek.\Dir{} spoon.\Ind{} \textsc{inst} eat-\Antip{}-\Pf{}//
\glft \trsl{Marek ate with a spoon.}//
\endgl
\a
\begingl
\gla *Marek travě piášček.//
\glb Marek.\Dir{} bread.\Ind{} eat-\Antip{}-\Pf{}//
\glft \trsl{(intended: Marek ate bread.)}//
\endgl
\a
\begingl
\gla *Marek piašték.//
\glb Marek.\Dir{} eat-\Pf{}//
\glft \trsl{(intended: Marek ate.)}//
\endgl
\xe

The instrumental phrase in the second example modifies the activity without
restoring a patient. The third example fails because indirect case cannot
license a patient excluded by the antipassive; the fourth lacks the overt
morphology required by the transitive base. A nominal dependent may separately
restrict an antipassive event under participant nominalisation, where a nominal
relational head licenses it (§\,\ref{sec:nominalised}). This nominal relation
does not extend to finite antipassive clauses.

\subsubsection{Minimalist analysis: the antipassive as inner Voice}
\label{sec:vm-antipassive}

Antipassive formation combines the transitive event predicate with an inner
Voice head that closes the patient role without projecting an independent
patient DP. For \ird{piašt-}, its interpretation can be represented as
\[
 \lambda x.\lambda e\,[\operatorname{eat}(e)\land
 \operatorname{Agent}(e,x)\land
 \exists y\,\operatorname{Theme}(e,y)].
\]
The existentially bound variable supplies the conceptual participant of eating;
it is unavailable for independent nominal realisation in the finite clause.
The representation is an analysis of Iridian's strict patient-suppression rule.
The availability of an oblique patient varies across antipassive constructions
cross-linguistically, as does their association with ergativity
\cite{heaton2020}; neither the suppression rule nor its aspectual restriction
follows from the term \term{antipassive} alone.

Asp.Align licenses inner antipassive Voice only in the ERG-aspect domain.
This dependency holds of the completed configuration, in which the higher Asp
head selects a compatible lower predicate. It does not require the lower head
to inspect a patient already assigned case by a not-yet-merged Asp head. The
remaining agent receives direct case through the resulting intransitive frame.
Embedded Asp imposes the same selectional restriction under participant
nominalisation, although finite case valuation is inactive there.

Antipassive Voice and Appl are separate heads, whose co-occurrence Iridian
excludes within one derived predicate. This is a language-specific selectional
restriction, rather than a consequence of two exponents occupying one universal
Voice position. Cross-linguistic combinations of antipassives and applicatives
are discussed by Basilico \cite[ch.~5]{basilico2026}. A higher causative may nevertheless select an antipassive
predicate (§\,\ref{sec:vm-caus}); its new causer does not restore the
suppressed lexical patient.

\subsection{The applicative (\ird{-éb-})}

The applicative \ird{-éb-} licenses one additional participant in relation to
the event and, in a simple indicative applicative, selects that participant as
the direct-case pivot. Its productive default is benefactive: the event is
performed for an individual's benefit, without entailing that the individual
performs the lexical action or receives its theme. The actor receives indirect
case in the NOM-aspect domain and transitive case in the ERG-aspect domain;
a retained lexical theme receives indirect case in either domain.

The following constructed contrast retains the same actor, theme, and
beneficiary. In the first clause, \ird{my} heads a benefactive phrase; in the
second, the beneficiary is a core DP selected by \ird{-éb-}:

\pex\label{ex:val-appl-transitive}
\a
\begingl
\gla Marek travě bylam my piaštíme.//
\glb Marek.\Dir{} bread.\Ind{} child.\Ind{} for eat-\Prog{}//
\glft \trsl{Marek is eating bread for the child's benefit.}//
\endgl
\a
\begingl
\gla Byl Marcí travě piaštébime.//
\glb child.\Dir{} Marek.\Ind{} bread.\Ind{} eat-\Appl{}-\Prog{}//
\glft \trsl{Marek is eating bread for the child's benefit.}//
\endgl
\xe

The applied child alone bears direct case. Marek and the bread remain arguments
of the eating event, although their shared indirect case does not distinguish
actor from theme; the predicate and context supply those relations. Neither
clause means that the child eats. Causative \ird{piáštrime}, by contrast,
can describe Marek's bringing about the child's eating
(§\,\ref{sec:aspect-caus}). A conventional \trsl{feed} reading is not part of
the productive semantics of \ird{-éb-}.

Two independently direct arguments are excluded from this simple applicative
frame:

\ex
\begingl
\gla *Marek byl piaštébime.//
\glb Marek.\Dir{} child.\Dir{} eat-\Appl{}-\Prog{}//
\glft \trsl{(intended: Marek is eating for the child's benefit.)}//
\endgl
\xe

An intransitive activity permits the same benefactive derivation without a
lexical theme. The following constructed examples illustrate that difference
and the ERG-domain frame of the transitive applicative:

\pex\label{ex:val-appl-ergative}
\a
\begingl
\gla Byl Marcí zemacébime.//
\glb child.\Dir{} Marek.\Ind{} run-\Appl{}-\Prog{}//
\glft \trsl{Marek is running for the child's benefit.}//
\endgl
\a
\begingl
\gla Byl Marku travě piaštébik.//
\glb child.\Dir{} Marek.\Trs{} bread.\Ind{} eat-\Appl{}-\Pf{}//
\glft \trsl{Marek ate bread for the child's benefit.}//
\endgl
\xe

\ird{-éb-} also admits role interpretations constrained by the predicate:
a recipient requires transfer or availability of a theme, a goal requires
motion or placement, and an instrument requires an activity in which that
participant serves as an implement. Locative uses require a lexically or
constructionally licensed spatial relation; the location of an event and its
destination are distinct relations. Discourse salience can motivate selection
of such a construction but cannot supply an otherwise unlicensed semantic
role. These restrictions distinguish Iridian's interpretations within the
broader typological variation in applicative bases and roles
\cite[§§2, 4]{polinsky2013appl}.

The goal and instrumental extensions can be illustrated with controlled
constructed pairs. Directed-motion \ird{stoj-} combines with a goal phrase
in the first example below; applicativisation makes that same town the
direct-case participant in the second. Both retain contemplative aspect:

\pex\label{ex:val-appl-goal}
\a
\begingl
\gla Marek hraďe de stojáž.//
\glb Marek.\Dir{} town.\Ind{} to go-\Ctpl{}//
\glft \trsl{Marek intends to go to the town.}//
\endgl
\a
\begingl
\gla Hrad Marcí stojébaž.//
\glb town.\Dir{} Marek.\Ind{} go-\Appl{}-\Ctpl{}//
\glft \trsl{Marek intends to go to the town.}//
\endgl
\xe

The town is the destination of projected motion, without an implication that
it benefits from that motion. Instrumental \ird{se}, already illustrated
with the antipassive, instead relates an implement to an activity. In the
following constructed pair \ird{rán}/\ird{ráne} supplies a controlled
case alternation (§\,\ref{sec:nm-stem-grades}); the context identifies the
hand as the means by which Marek eats:

\pex\label{ex:val-appl-instrument}
\a
\begingl
\gla Marek travě ráne se piaštíme.//
\glb Marek.\Dir{} bread.\Ind{} hand.\Ind{} \textsc{inst} eat-\Prog{}//
\glft \trsl{Marek is eating bread by hand.}//
\endgl
\a
\begingl
\gla Rán Marcí travě piaštébime.//
\glb hand.\Dir{} Marek.\Ind{} bread.\Ind{} eat-\Appl{}-\Prog{}//
\glft \trsl{Marek is eating bread by hand.}//
\endgl
\xe

Neither extension changes the lexical event: the hand does not become an
eater, and the town does not become a moving participant. The instrumental
interpretation depends on an implement-compatible event, just as the goal
interpretation depends on directed motion; unrestricted exchange of these
roles is excluded by those selectional conditions. Recipient and locative
extensions remain restricted to compatible predicates, whose full lexical
distribution is unspecified.

Directive Mood selects the lexical theme where one is retained, so that an
applied beneficiary receives indirect case (§\,\ref{sec:mood-directive-valency}):

\pex
\a
\begingl
\gla Trava bylam piaštébím.//
\glb bread.\Dir{} child.\Ind{} eat-\Appl{}-\Imp{}//
\glft \trsl{Eat bread for the child's benefit.}//
\endgl
\a
\begingl
\gla Trava bylam piaštébká.//
\glb bread.\Dir{} child.\Ind{} eat-\Appl{}-\Hort{}//
\glft \trsl{Let's eat bread for the child's benefit.}//
\endgl
\xe

The beneficiary is required overtly in these directive constructions. Its
indirect case reflects the higher Mood-conditioned selection rule, although
Appl continues to introduce it as an argument.

\subsubsection{Minimalist analysis: the applicative as ApplP}
\label{sec:vm-appl}

The productive benefactive is analysed as a high applicative, whose added
participant is related to an event. Pylkkänen distinguishes this relation from
low applicatives, which relate individuals through transfer of possession
\cite[ch.~2]{pylkkanen2008}. Iridian's benefactive introduces no possession
entailment, and the grammar assigns its productive interpretation to a single
Appl head. Recipient readings are predicate-conditioned extensions rather than
evidence for an independently productive low-applicative position.

A schematic interpretation is
\[
 \mathrm{Appl}_{\mathrm{BEN}}
 =\lambda P.\lambda y.\lambda e\,
 [P(e)\land\operatorname{Benefit}(e,y)],
\]
where $P$ is an event predicate and $y$ is the beneficiary. Appl combines with
the event predicate below the layer introducing its actor; both lie below an
optional Caus head and the outer orientation head. The constituent order is
head-final throughout:

\begin{center}
\begin{forest}
for tree={align=center}
[\synhead{ApplP}
  [{\synhead{DP}\\{\footnotesize beneficiary}}]
  [{Appl$'$}
    [\synhead{VP} [{\ldots}]]
    [{Appl\\\ird{-éb-}}]
  ]
]
\end{forest}
\end{center}

Appl introduces the beneficiary; the completed clause's licensing configuration
determines case and pivot privileges. A simple indicative applicative selects
the beneficiary in either alignment domain, whereas higher Caus or directive
Mood can supersede that selection. The outer orientation head is zero in a simple applicative, where its
selection follows Appl; \ird{-éb-} selects perfective \ird{-ik} and
licenses the ERG-domain frame without an additional \ird{-n-}. A higher
causative creates its own orientation domain and requires overt Voice
\ird{-an-} in the ERG frame. Iridian excludes a
second applicative and the combination of applicative and antipassive within
the same predicate.

\subsection{The rhotic causative}\label{sec:aspect-caus}

The rhotic causative adds a causer while retaining the base's independently
licensed participants. Compact \ird{-r-} and expanded \ird{-ér-}/\ird{-er-}
differ in their local distribution of strength, while a lexically restricted
expanded class also selects front-grade PAL. These valency increases have a
comparison in Dixon's morphological causative type \cite{dixon2000};
the allomorphs, quantities, and case rules below specify Iridian.

\subsubsection{Causative grades and outer orientation}
\label{sec:causative-grades}

Allomorph selection precedes quantity assignment. Compact \ird{-r-}
selects G1 on an eligible local stem nucleus, giving \ird{piáštr-},
\ird{vádr-}, and \ird{zemácr-}. It does not change the preceding
thematic consonant: \ird{piáštrime} retains /t/, and
\ird{zemácrime} retains /ts/. The compact onset clusters /tr/ and /br/
follow the ordinary syllabification rules
(§\,\ref{sec:cluster-repair}).

Expanded \ird{-er-} occurs where the compact sequence cannot be
syllabified and in specified lexical classes, including \ird{ščen-},
stative stems in \ird{-n-}, \ird{stoj-}, \ird{nočk-}, and the
replacement antipassive \ird{piášč-}. With a short input stem, the
expanded exponent carries strength as \ird{-ér-}:
\ird{ščenérime}, \ird{svornérime}. Already strong \ird{piášč-}
instead takes short \ird{-er-}, as in \ird{piáščeranik}.
The listed compact alternative \ird{ščénrime} retains G1 on the root;
its existence does not make every expanded-class stem freely compact.

Before consonant-initial directive Mood, the expanded allomorph likewise
supplies the causative vowel: \ird{piaštérká}, \ird{piaštérku}.
These forms select expanded Caus before assigning strength to it.
Compact \ird{piáštrím} and \ird{piáštré}, by contrast, retain
their root G1 alongside the independently long Mood vowel. The distribution
does not follow a ban on two long vowels in a word.

The mutating expanded class is lexical. Its front-grade consonantal
specification is independent of the surface vowel quantity:
\ird{vad-} with that causative gives \ird{vaďérime}, beside compact
\ird{vádrime}. PAL supplies /ɟ/, written \ird{ď}; G1 belongs to
the suffix vowel, so the root remains short. Long \ird{é} replaces
the short-vowel spelling \ird{ě} in the strong expanded exponent.
The front and plain expanded classes can therefore share the spelling
\ird{-ér-}/\ird{-er-} while selecting different consonantal
correspondences. No spelling-based rule freely assigns PAL to another
causative class or changes /r/ to learned \ird{ř}.

An intervening valency head establishes a separate domain. Protected
applicative \ird{-éb-} retains its quantity in \ird{piaštébrime}
and prevents outer Caus from lengthening the original root. The same
locality prevents a causative outside antipassive morphology from applying
PAL or assigning a new root grade through that inner head.
Morphology-sensitive palatalisation has a comparative parallel in the Polish
suffix classes discussed by Rubach \cite[§§2--3]{rubach2019}; the lexical
classes and local quantitative specifications stated here belong to Iridian.

Outer patient-selecting Voice takes \ird{-an-} immediately after any
rhotic causative, wherever that orientation head is overt, and ordinary
\ird{-n-} elsewhere. Thus \ird{piášt-r-an-ik} yields
\ird{piáštranik}, while \ird{piášt-r-an-ime} yields passive
\ird{piáštranime}. The short aspect endings preserve Caus-selected
strength. Expanded \ird{piášč-er-an-ik} gives \ird{piáščeranik},
and front-grade expanded \ird{vaď-ér-an-ik} gives \ird{vaďéranik}.
The causative and orientation remain distinct exponents, with no metathesis
or joint causative--Voice morpheme.

In the Distributed Morphology analysis, these causative and orientation
allomorphs are Vocabulary Items with lexical or morphological conditioning,
following Embick and Noyer's account of contextual allomorphy
\cite[§2.3]{embicknoyer2007}. Their lexical membership does not follow from
a root's meaning.

\begin{table}[htbp]
\footnotesize
\caption{Constructed causative forms of compact-class \ird{piašt-}.}
\label{tab:rhotic-causative-paradigm}
\begin{tabularx}{\textwidth}{L L L L}
\toprule
Category & Form & Category & Form \\
\midrule
Continuous & \ird{piáštra} & Progressive & \ird{piáštrime} \\
Contemplative & \ird{piáštraž} & Perfective & \ird{piáštranik} \\
Retrospective & \ird{piáštranav} & Imperative & \ird{piáštrím} \\
Prohibitive & \ird{piáštré} & Hortative & \ird{piaštérká} \\
PFV imperative & \ird{piáštrovím} & Sequential & \ird{piáštru}, \ird{piáštrení} \\
\bottomrule
\end{tabularx}
\end{table}

\subsubsection{Causative argument frames}

In an active NOM-domain causative, the causer is direct and all other arguments
are indirect. In an ERG-domain causative, the causer is transitive and one
embedded participant is direct: the retained lexical theme takes precedence,
followed by an applied participant if there is no lexical theme, and otherwise
the causee. Any remaining causee or applied participant is indirect. This
hierarchy belongs to the causative frame, so an embedded applicative's simple
beneficiary selection does not independently assign a second direct case.
Here \term{retained theme} designates the separate patient of a transitive
base. An intransitive's sole participant becomes the causee even when its
lexical role is theme-like, as with \ird{ščen-} \trsl{arrive}; that role
does not create an additional patient ahead of an applied participant in the
causative hierarchy.

\begin{table}[htbp]
\footnotesize
\caption{Active indicative causative frames. A dash indicates an absent role.}
\label{tab:causative-frames}
\begin{tabularx}{\textwidth}{L C C C C}
\toprule
Base and domain & Causer & Causee & Theme & Applied \\
\midrule
Intransitive, NOM & \textsc{dir} & \textsc{ind} & --- & --- \\
Intransitive, ERG & \textsc{trs} & \textsc{dir} & --- & --- \\
Transitive, NOM & \textsc{dir} & \textsc{ind} & \textsc{ind} & --- \\
Transitive, ERG & \textsc{trs} & \textsc{ind} & \textsc{dir} & --- \\
Applied transitive, NOM & \textsc{dir} & \textsc{ind} & \textsc{ind} & \textsc{ind} \\
Applied transitive, ERG & \textsc{trs} & \textsc{ind} & \textsc{dir} & \textsc{ind} \\
Applied intransitive, NOM & \textsc{dir} & \textsc{ind} & --- & \textsc{ind} \\
Applied intransitive, ERG & \textsc{trs} & \textsc{ind} & --- & \textsc{dir} \\
Antipassive, ERG & \textsc{trs} & \textsc{dir} & --- & --- \\
\bottomrule
\end{tabularx}
\end{table}

The base contrasts can be compared directly in the following constructed
paradigm. \ird{Zemac-} denotes an actor's activity, whereas \ird{ščen-}
denotes its sole participant's arrival; their ordinary case frames coincide.
The stative \ird{svorn-} independently selects perfective \ird{-ik}.
The transitive base alone requires \ird{-n-} in the ERG frame.

\begin{table}[htbp]
\footnotesize
\caption{Constructed ordinary base clauses in the two alignment domains.}
\label{tab:valency-base-contrasts}
\begin{tabularx}{\textwidth}{l L L}
\toprule
Base & NOM: progressive & ERG: perfective \\
\midrule
\ird{piašt-} & \ird{Marek travě piaštíme.} & \ird{Trava Marku piaštník.} \\
\ird{zemac-} & \ird{Marek zemácime.} & \ird{Marek zemácek.} \\
\ird{ščen-} & \ird{Byl ščénime.} & \ird{Byl ščének.} \\
\ird{svorn-} & \ird{Problem svórnime.} & \ird{Problem svórnik.} \\
\bottomrule
\end{tabularx}
\end{table}

In the transitive row the agent is direct in NOM and transitive in ERG,
while the theme is indirect in NOM and direct in ERG. Each intransitive row
has one core participant, which is direct and privileged in both domains.
The progressive stative describes a developing increase in size, while its
perfective presents a bounded becoming-big event; stative causativisation adds
a distinct causer (§\,\ref{sec:stative-valency}).

Complex forms below retain visible morphological boundaries where these
clarify suffix order; the hyphens are analytical notation. The following
constructed pairs make the causee and any retained theme overt:

\pex\label{ex:val-caus-intransitive}
\a
\begingl
\gla Marek bylam zemácrime.//
\glb Marek.\Dir{} child.\Ind{} run-\Caus{}-\Prog{}//
\glft \trsl{Marek is making the child run.}//
\endgl
\a
\begingl
\gla Byl Marku zemác-r-an-ik.//
\glb child.\Dir{} Marek.\Trs{} run-\Caus{}-\textsc{trs}-\Pf{}//
\glft \trsl{Marek made the child run.}//
\endgl
\xe

The activity predicate \ird{zemac-} \trsl{run} supplies an actor causee.
By comparison, the sole participant of \ird{ščen-} \trsl{arrive} undergoes a
change of location. The grammar treats the former as unergative and the latter
as unaccusative, using that distinction to characterise the base relation;
both follow the same causative case rule. The following constructed paradigm
includes the underived ERG-domain intransitive, which needs no transitive
\ird{-n-}, and causatives in both alignment domains:

\pex\label{ex:val-unaccusative}
\a
\begingl
\gla Byl ščének.//
\glb child.\Dir{} arrive-\Pf{}//
\glft \trsl{The child arrived.}//
\endgl
\a
\begingl
\gla Marek bylam ščenérime.//
\glb Marek.\Dir{} child.\Ind{} arrive-\Caus{}-\Prog{}//
\glft \trsl{Marek is causing the child to arrive.}//
\endgl
\a
\begingl
\gla Byl Marku ščen-ér-an-ik.//
\glb child.\Dir{} Marek.\Trs{} arrive-\Caus{}-\textsc{trs}-\Pf{}//
\glft \trsl{Marek caused the child to arrive.}//
\endgl
\xe

\pex\label{ex:val-caus-transitive}
\a
\begingl
\gla Marek bylam travě piáštrime.//
\glb Marek.\Dir{} child.\Ind{} bread.\Ind{} eat-\Caus{}-\Prog{}//
\glft \trsl{Marek is making the child eat bread.}//
\endgl
\a
\begingl
\gla Travó Marku bylam piášt-r-an-ik.//
\glb bread.\Dir{}.\Def{} Marek.\Trs{} child.\Ind{} eat-\Caus{}-\textsc{trs}-\Pf{}//
\glft \trsl{Marek made the child eat the bread.}//
\endgl
\xe

The rhotic causative expresses causation without specifying coercion, permission, or the
causer's degree of intervention. The causee's agency, the lexical event, and
context support more particular readings; viewpoint aspect does not entail
causal directness or control. Thus the perfective transitive example can
receive a permissive interpretation in a context where Marek enables the
child to eat, while retrospective \ird{piášt-r-an-av} can describe direct
compulsion when that relation is supplied by context. Both preserve the same
causative argument frame.

An applicative precedes the causative. The following constructed clauses use
plural \ird{býlac} for the causee and singular \ird{bylam} for the beneficiary,
so that all four participants are represented:

\pex\label{ex:val-caus-appl}
\a
\begingl
\gla Marek býlac bylam travě piaštébrime.//
\glb Marek.\Dir{} child.\Ind{}.\Pl{} child.\Ind{} bread.\Ind{} eat-\Appl{}-\Caus{}-\Prog{}//
\glft \trsl{Marek is making the children eat bread for the child's benefit.}//
\endgl
\a
\begingl
\gla Trava Marku býlac bylam piaštéb-r-an-ik.//
\glb bread.\Dir{} Marek.\Trs{} child.\Ind{}.\Pl{} child.\Ind{} eat-\Appl{}-\Caus{}-\textsc{trs}-\Pf{}//
\glft \trsl{Marek made the children eat bread for the child's benefit.}//
\endgl
\xe

Their scope is $\mathrm{CAUS}[\mathrm{APPL}[\mathrm{EAT}]]$: the child benefits
from the embedded eating event. A beneficiary of the causing event itself can
be expressed separately with \ird{my}; the suffix order does not license an
outer applicative reading. Multiple indirect-case arguments need not occur in
all uses, but an omitted causee requires the ordinary contextual licensing of
argument omission. The morphology itself does not introduce an arbitrary
silent \trsl{someone}.

A causative of an applied intransitive lacks a separate lexical patient.
The following constructed pair therefore selects the causer in the
progressive and the beneficiary in the perfective, leaving the running
children in indirect case in both:

\pex\label{ex:val-caus-appl-intransitive}
\a
\begingl
\gla Marek býlac bylam zemacébrime.//
\glb Marek.\Dir{} child.\Ind{}.\Pl{} child.\Ind{} run-\Appl{}-\Caus{}-\Prog{}//
\glft \trsl{Marek is making the children run for the child's benefit.}//
\endgl
\a
\begingl
\gla Byl Marku býlac zemacéb-r-an-ik.//
\glb child.\Dir{} Marek.\Trs{} child.\Ind{}.\Pl{} run-\Appl{}-\Caus{}-\textsc{trs}-\Pf{}//
\glft \trsl{Marek made the children run for the child's benefit.}//
\endgl
\xe

The perfective child is the beneficiary, while plural \ird{býlac}
expresses the causee. The contrast with the preceding transitive pair
follows from the presence or absence of a retained patient, rather than from
a change in the benefactive relation contributed by Appl.

An antipassive base retains its actor as the causee while its lexical patient
remains unavailable. Causativisation therefore produces a new two-participant
predicate without reversing the lower reduction:

\ex\label{ex:val-caus-antipassive}
\begingl
\gla Byl Marku piáščeranik.//
\glb child.\Dir{} Marek.\Trs{} eat-\Antip{}-\Caus{}-\textsc{trs}-\Pf{}//
\glft \trsl{Marek made the child eat.}//
\endgl
\xe

The ERG-domain licensing restriction on antipassive remains operative in this
configuration; a NOM-domain counterpart containing \ird{-aš-} is excluded.
The direct-case child is the causee, never a restored lexical patient.

\subsubsection{Minimalist analysis: the causative as CausP}
\label{sec:vm-caus}

Caus combines with an already formed event predicate, introducing its causer
above the base actor and any applied argument. The outer orientation head
subsequently selects the derived predicate's privileged participant, while
Asp.Align supplies the alignment domain. The distinction between productive
causative structure and other verbal heads follows the analytical issues
surveyed by Harley \cite[§1.4]{harley2017}; assigning causer introduction
to Caus and patient selection to higher Voice is the Iridian analysis.
The head-final hierarchy yields
root--inner derivation--causative--orientation--aspect:

\begin{center}
\begin{forest}
for tree={align=center}
[\synhead{VoiceP}
  [\synhead{CausP}
    [{\synhead{DP}\\{\footnotesize causer}}]
    [{Caus$'$}
      [{event predicate\\{\footnotesize including causee}}]
      [{Caus\\\ird{-r-} / \ird{-er-} / \ird{-ěr-}}]
    ]
  ]
  [{Voice\\\ird{-an-} / $\varnothing$}]
]
\end{forest}
\end{center}

The overt outer \ird{-an-} licenses the selected embedded participant in an
ERG causative. In the NOM domain, zero outer orientation gives the active
causer-oriented frame; overt \ird{-an-} selects the embedded theme, then an
applied participant if no theme is retained, and otherwise the causee. It
licenses the causer as an optional indirect agentive dependent. This is the passive of the
derived causative. An embedded causee remains an argument of its own predicate;
passivisation of the higher causing event does not suppress that relation.

\subsubsection{The Voice head \ird{-n-}}
\label{sec:vm-n}

\ird{-n-} realises outer patient-selecting Voice, following causative
\ird{-r-}, \ird{-er-}, or \ird{-ěr-} wherever both occur,
with the conditioned Voice allomorph \ird{-an-}. With a simple transitive base it
selects the
lexical patient; with a causative it selects according to the hierarchy in
Table~\ref{tab:causative-frames}. Its interpretation as patient selection thus
extends to a causee that is the affected participant of an intransitive-base
causative. This higher orientation head is distinct from inner antipassive
Voice and from argument-introducing Appl and Caus.

In the ERG-aspect domain, \ird{-n-} licenses the ordinary two-participant
transitive frame: the selected participant is direct and the highest actor is
transitive. Simple applicatives have their own overt licensing through
\ird{-éb-}; simple antipassives instead retain only an actor, which receives
direct case. The contrast between overt \ird{-n-} and a bare perfective on a
transitive base is therefore a restriction on the completed licensing frame,
not evidence that patient argumenthood universally requires an overt suffix.

In the NOM-aspect domain, \ird{-n-} overrides active actor selection and
licenses a different external-argument configuration: the patient is direct,
while the highest actor can occur as an optional indirect agentive dependent.
The grammar calls this configuration \term{passive}. Patient selection and
external-argument licensing are separate parts of its specification; assigning
direct case to a patient alone would not establish the latter property.
The relative height of Voice accounts for \ird{piášt-r-an-ik} without
metathesis of Voice and Caus.

\subsection{The passive}\label{sec:vm-passive}

The NOM-domain passive combines patient-selecting \ird{-n-}
(\ird{-an-} after a rhotic causative) with continuous
\ird{-a}, progressive \ird{-ime}, or contemplative \ird{-áž}/\ird{-až}, selected by the local
quantity specification. Its patient
receives direct case, while an expressed actor receives indirect case as an
agentive dependent licensed by the passive construction. An unexpressed actor
can be contextually recovered or left unspecified. The aspectual domain remains
NOM, although active actor selection is overridden.

The indirect agent differs in licensing from the ordinary indirect patient of
an active NOM-domain transitive: the latter is an internal argument of the
lexical predicate, whereas the former is an optional external dependent of the
passive configuration. Their shared case ending does not establish identical
argument status. Further diagnostics of the passive agent's binding and control
properties are not established by the case alternation alone.

\pex
\a
\begingl
\gla Travó Marcí piaštníme.//
\glb bread.\Dir{}.\Def{} Marek.\Ind{} eat-\textsc{pass}-\Prog{}//
\glft \trsl{The bread is being eaten by Marek.}//
\endgl
\a
\begingl
\gla Travó piaštníme.//
\glb bread.\Dir{}.\Def{} eat-\textsc{pass}-\Prog{}//
\glft \trsl{The bread is being eaten.}//
\endgl
\xe

\ird{Marcí} is the indirect form of \ird{Marek}; transitive \ird{Marku}
belongs to the ERG-domain active frame. The same outer head follows the
causative in the constructed passive below, where Marek is the optional causer
dependent and the child remains the embedded actor:

\ex
\begingl
\gla Trava Marcí bylam piášt-r-an-ime.//
\glb bread.\Dir{} Marek.\Ind{} child.\Ind{} eat-\Caus{}-\textsc{pass}-\Prog{}//
\glft \trsl{Bread is being caused by Marek to be eaten by the child.}//
\endgl
\xe

The passive privileges the bread for
participant nominalisation and direct-gap relativisation. It can thereby
support continued reference to that participant, although aboutness topicality
requires its independent discourse conditions.


\subsection{Pragmatic use of valency operations}

The antipassive presents an activity without an independently licensed patient,
which can support discourse centred on the actor or on the activity itself.
An ordinary transitive with an omitted patient differs because its patient
remains a referential argument recoverable from context. Choosing between these
constructions therefore concerns argument structure as well as discourse.

An applicative makes an event-related participant available for pivot-restricted
constructions, including participant nominalisation and relative formation.
Continued reference to a beneficiary may favour this choice, although a topic
relation is neither required nor supplied by the suffix. The beneficiary
interpretation comes from the applicative construction; discourse salience
does not by itself license an instrumental, goal or stimulus reading.

A causative introduces a causing event and a causer, while the lexical predicate
and context determine whether the causal relation involves inducement,
permission, delegated agency or physical intervention. Perfective and
retrospective viewpoints can present the same relation from different temporal
perspectives without encoding different degrees of force. Directive Mood
likewise changes the construction's orientation without changing the causal
relation (§\,\ref{sec:mood-directive-valency}).

\section{Negation}\label{sec:negation}

\subsection{Forms, distribution, and prosody}\label{sec:neg-allomorphy}

Indicative verbal negation is expressed by unstressed \ird{za-}, which attaches
as a proclitic to the left edge of the verbal complex without displacing the
host's primary stress. The imperative, hortative, and irrealis have the negative
mood suffixes \ird{-é}, \ird{-ku}, and \ird{-ir}, respectively
(§\,\ref{sec:mood}). Nominal predicates take clause-final \ird{staly}, while
existential predicates distinguish affirmative \ird{jec-} from negative
\ird{zad-} (§\,\ref{sec:clausal-negation} of Chapter~\ref{ch:simple}).

The negative marker is underlyingly /za/, with its short vowel retained before
a consonant, as in \ird{za-} + \ird{piaštník} $\rightarrow$
\ird{zapiaštník}, and deleted before a vowel, giving \ird{z-}. Deletion applies before both short and long vowels, irrespective of stress,
and is morphologically restricted to the negative marker. The choice of form consequently depends on the host's
phonological onset, including the distinction between vocalic \ird{i} and an
orthographic palatalisation marker (§\,\ref{sec:orthographic-i}).

\begin{table}[htbp]
  \centering
  \caption{Indicative negative allomorphy; vowel and glide sequences are schematic.}\label{tab:neg-allomorphy}
  \begin{tabularx}{\textwidth}{l l X}
    \toprule
    {\sc host onset} & {\sc spelling} & {\sc analysis} \\
    \midrule
    Consonant & \ird{za-} & Short /a/ retained; e.g.\ \ird{zapiaštník}. \\
    \ird{a, á} & \ird{za\ldots, zá\ldots} & \ird{z-} + host vowel; length belongs to the host. \\
    \ird{e, é} & \ird{ze\ldots, zé\ldots} & \ird{z-} + unchanged host vowel. \\
    \ird{i, í} & \ird{zi\ldots, zí\ldots} & \ird{z-} + unchanged host vowel; /z/ remains alveolar. \\
    \ird{o, ó} & \ird{zo\ldots, zó\ldots} & \ird{z-} + unchanged host vowel. \\
    \ird{u, ú} & \ird{zu\ldots, zú\ldots} & \ird{z-} + unchanged host vowel. \\
    \ird{y, ý} & \ird{zy\ldots, zý\ldots} & \ird{z-} + unchanged host vowel. \\
    \ird{j} /j/ & \ird{zaj\ldots} & Consonantal onset; full \ird{za-} retained. \\
    \bottomrule
  \end{tabularx}
\end{table}

Before \ird{i, í}, the reduced marker retains alveolar /z/: the high front
vowel does not select a postalveolar consonant
(§\,\ref{sec:sibilant-palatalisation}). The glide /j/, being consonantal,
takes \ird{za-}, so that a \ird{j}-initial host begins \ird{zaj\ldots}
under negation. The negative marker has no \ird{ž-} allomorph.

In weak position, the marker's short /a/ may centralise to [ɐ], giving
[za-] $\sim$ [zɐ-] without any change to the host's stress
(§\,\ref{sec:stress}). Since the marker has no long-vowel variant, an initial
\ird{zá\ldots} contains \ird{z-} followed by the host's own long \ird{á}.
Throughout this grammar, the label \ird{za-} includes its regular \ird{z-}
allomorph.

\subsection{Formal analysis: \synhead{NegP} and the position of the negative marker}
\label{sec:vm-neg}
\label{sec:vm-neg-exponents}

Indicative clausal negation is analysed as a \synhead{NegP} above
\synhead{AspP} and below \synhead{MoodP}
(Table~\ref{tab:th-clause-spine}), whose realisation depends on the predicate.
In ordinary verbal clauses, \ird{za-} precedes the root and forms a prosodic
unit with the inflected verb, while in zero-copula clauses \ird{staly} follows
the nominal predicate. With existential predicates, negative polarity is
expressed by suppletion of the root, \ird{jec-}/\ird{zad-}.

Iridian exhibits \term{strict negative concord}: negative indefinites such as
\ird{deno} \trsl{nothing} and \ird{zide} \trsl{no one} require the
negative realisation appropriate to their finite clause, while their
co-occurrence expresses a single semantic negation
(§\,\ref{sec:nm-negative-concord}). The grammar analyses this dependency through
a silent operator $\mathrm{Op}_{\neg}$ bearing an interpretable negative
feature $[i\textsc{neg}]$, together with an uninterpretable feature
$[u\textsc{neg}]$ on Neg\textsuperscript{0} and on each concord associate.
The operator occupies Spec,\synhead{NegP}; it supplies semantic negation,
whereas the feature on Neg\textsuperscript{0} determines its participation in
the dependency and its morphological realisation. This division of labour
follows the syntactic approach to negative dependencies developed by
Zeijlstra and Penka \cite{zeijlstra2022,penka2011}; its application to Iridian is the analysis
adopted here.

The preverbal position of \ird{za-} is attributed to post-syntactic attachment
of the exponent of Neg\textsuperscript{0} to the left edge of the verbal host,
within the framework of Distributed Morphology \cite{embicknoyer2007}.
Although Neg is structurally higher than Asp, prefixal realisation is
compatible with the Mirror Principle \cite{baker1985}, which constrains the
relative order of affixes without requiring a higher head to be suffixal
\cite{harley2011}.

In zero-copula clauses such as \ird{Marek doktor staly} \trsl{Marek is not a
doctor}, invariant \ird{staly} realises a head-final Neg selecting the nominal
predication directly, without an intervening \synhead{AspP}. In verbal
clauses, Neg instead selects an aspectual predicate, including a finite
stative predicate (§\,\ref{sec:vm-tension-stative}). The difference in
complement type corresponds to the distribution of \ird{staly} and
\ird{za-}.

Existential negation is expressed by \ird{zad-}, the negative allomorph of
Exist\textsuperscript{0}, with a zero exponent at Neg\textsuperscript{0}.
Thus \ird{za-} and \ird{staly} spell out Neg itself, whereas \ird{zad-}
realises the existential root in the environment of Neg. Root suppletion is
well documented among negative existentials \cite{veselinova2016}; in the
Iridian analysis, the negative head conditions the suppletive root while
remaining distinct from it in the syntactic representation.

\begin{figure}[ht]
  \footnotesize
  \caption{Proposed structures for indicative verbal, equative, and existential
  negation. The verbal marker attaches to its host at PF, while existential
  negation combines a zero exponent of Neg with a negative root allomorph.
  Higher topic structure is omitted; Spec,\synhead{NegP} hosts the silent
  negative operator.}\label{fig:neg0-comparative}
  \medskip
  \centering
  \scriptsize
  \begin{minipage}[t]{0.30\textwidth}
    \centering
    \textsc{Verbal clause}\par
    \medskip
    \begin{forest}
      syntax tree, for tree={l sep=8pt,s sep=6pt,inner sep=0.5pt}
      [\synhead{NegP}
        [{$\mathrm{Op}_{\neg}$\\{$[i\textsc{neg}]$}}]
        [{Neg$'$}
          [\synhead{AspP}
            [{Asp$'$}
              [\synhead{VP}
                [{V\\{\footnotesize\textit{root}}}]
              ]
              [{Asp\\{\footnotesize\textit{asp.align}}}]
            ]
          ]
          [{Neg$^{0}$\\{$[u\textsc{neg}]$}\\\ird{za-}\\{\small(PF prefix)}}]
        ]
      ]
    \end{forest}
  \end{minipage}
  \hfill
  \begin{minipage}[t]{0.30\textwidth}
    \centering
    \textsc{Equative clause}\par
    \medskip
    \begin{forest}
      syntax tree, for tree={l sep=8pt,s sep=6pt,inner sep=0.5pt}
      [\synhead{NegP}
        [{$\mathrm{Op}_{\neg}$\\{$[i\textsc{neg}]$}}]
        [{Neg$'$}
          [\synhead{EquP}
            [{\synhead{DP}\\{\footnotesize lower copy}}]
            [{Eq$'$}
              [{\synhead{DP}\\{\footnotesize pred.}}]
              [{Eq$^{0}$\\{$\varnothing$}}]
            ]
          ]
          [{Neg$^{0}$\\{$[u\textsc{neg}]$}\\\ird{staly}}]
        ]
      ]
    \end{forest}
  \end{minipage}
  \hfill
  \begin{minipage}[t]{0.35\textwidth}
    \centering
    \textsc{Existential clause}\par
    \medskip
    \begin{forest}
      syntax tree, for tree={l sep=8pt,s sep=6pt,inner sep=0.5pt}
      [\synhead{NegP}
        [{$\mathrm{Op}_{\neg}$\\{$[i\textsc{neg}]$}}]
        [{Neg$'$}
          [\synhead{AspP}
            [{Asp$'$}
              [ExistP
                [{Exist$'$}
                  [{Exist$^{0}$\\\ird{zad-}\\{\small(suppletive)}}]
                ]
              ]
              [{Asp\\{\footnotesize\textit{asp}}}]
            ]
          ]
          [{Neg$^{0}$\\{$[u\textsc{neg}]$}\\{\small($\varnothing$)}}]
        ]
      ]
    \end{forest}
  \end{minipage}
\end{figure}

In Figure~\ref{fig:neg0-comparative}, Neg follows its complement in all
three structures, with the prefixal position of \ird{za-} determined at PF.
An ordinary linked topic occupies the independently projected
 Spec,\synhead{TopP}, above the structures shown here
 (§\,\ref{sec:topic-comment}). Quantificational pivots remain in the lower
 argument domain, with negative indefinites interpreted within the scope of
 the negative operator; topicless existential pivots likewise remain below
 that operator (§\,\ref{sec:existential-constructions}). The nominal
 constituents of an equative may permute with topic and comment, while
 \ird{staly} remains clause-final.

In negative existential clauses, \ird{zad-} and the zero exponent of Neg form
a single pattern of contextual exponence, excluding both \ird{*za-jec-} and
\ird{*za-zad-}. Negative mood clauses have a different structure, in which
negative Mood selects Asp directly and no separate Neg head is merged
(§\,\ref{sec:vm-moodp}).

The clitic analysis predicts adjacency between \ird{za-} and its verbal host,
including the exclusion of intervening particles; its behaviour under ellipsis
is undocumented. The invariance of \ird{staly}, which bears no aspect or mood
inflection, supports its treatment as a functional negative marker. The
prefixal and independent forms belong to two of the positional types
represented in Dryer's survey of negative markers \cite{dryer2013neg}.

\subsubsection{Identificational negation and the scope of \ird{staly}}
\label{sec:vm-ident-neg}

When a participant nominal serves as the predicate of a zero-copula clause,
\ird{staly} negates an identification rather than a finite verbal assertion
(§\,\ref{sec:nmlz-predicative}). Whereas \ird{Marek doktor staly}
\trsl{Marek is not a doctor} denies a property ascription,
\ird{Lobek Jancí piaštníkou staly} \trsl{The apple is not what Janek ate}
denies that the apple satisfies the participant description under discussion.
Both are propositional negations, although the propositions differ in the kind
of predication they contain.

When Janek's eating and the existence of a patient are already established in
the discourse, \ird{Lobek Jancí piaštníkou staly} rejects the apple as that
patient, suggesting that he ate something else. The alternative-participant
reading depends on the continued availability of that existential background.
Existence and uniqueness requirements in contexts with no established event or
with several possible participants are undocumented.

\medskip

\noindent\begin{tabularx}{\linewidth}{l L L}
\toprule
& \textsc{finite verbal negative} & \textsc{identificational negative} \\
\midrule
Predicate & finite event or state predicate & participant nominal \\
What is denied & verbal predication & identification of the topic \\
Event status & need not be given & commonly given in corrective uses \\
Negative form & \ird{za-} on the verbal complex & clause-final \ird{staly} \\
\bottomrule
\end{tabularx}

\medskip

Finite verbal negation denies an event or state predication, whereas nominal
negation denies a property ascription or identification. In identificational
uses, the discourse background supplies the event against which the topic's
participant role is evaluated (§\,\ref{sec:ident-equative} of
Chapter~\ref{ch:simple}).

\subsubsection{Negative concord and semantic interpretation}
\label{sec:vm-negative-concord}

A negative indefinite introduces an existential variable whose restriction is
supplied by its lexical content, while the silent operator supplies the single
negation under which that variable is interpreted. Neg\textsuperscript{0}
and the indefinites bear $[u\textsc{neg}]$ and are jointly licensed by
$\mathrm{Op}_{\neg}[i\textsc{neg}]$ through Multiple Agree. The feature
notation therefore distinguishes the source of semantic negation from its
several morphological reflexes; neither an additional indefinite nor an
overt negative exponent introduces another operator. A schematic sequence
\ird{zide \ldots\ deno \ldots\ za-V}, for instance, contains two
concord associates within one negative clause and receives a single-negation
interpretation. The indefinite series and the contrasts with affirmative
clauses are described in §\,\ref{sec:nm-negative-concord}.

Vocabulary Insertion realises Neg\textsuperscript{0}$[u\textsc{neg}]$ as
\ird{za-} $\sim$ \ird{z-} with verbal predicates, \ird{staly} with
zero-copula equatives, and $\varnothing$ with existential predicates; in the
last environment, Exist\textsuperscript{0} is realised as \ird{zad-}.
The operator and the feature dependency belong to syntax and interpretation,
whereas these exponents and the attachment of \ird{za-} belong to the
morphological realisation of that structure. The analysis consequently retains
both the existing verbal template and the distinction between a zero exponent
of Neg and the absence of NegP.

Concord is local to the finite clause: an operator in a matrix clause does
not license a negative indefinite in a finite complement; conditionality alone likewise supplies no licenser
for an antecedent (Chapter~\ref{ch:complex}). Interrogative \ird{by} and non-assertive
complement marking do not supply $[i\textsc{neg}]$. A fronted concord associate
is licensed in its lower position and interpreted there under the negative
operator, even when its pronounced copy precedes the rest of the clause
(§\,\ref{sec:topic-comment}). Genuine double negation requires two semantic
operators, ordinarily in distinct clauses or propositions; the co-occurrence
of an indefinite and \ird{za-}, \ird{staly}, or \ird{zad-} within a concord
dependency cannot produce it.

\section{Mood}\label{sec:mood}

The mood system comprises the indicative, the imperative, the prohibitive, the
hortative, and the irrealis. Mood is expressed by a suffix occupying the
position immediately outside the aspect-alignment morpheme; in the indicative
this suffix is null. Where a non-indicative mood suffix is present, the
Asp.Align head remains in the clause, while Mood licenses a two-way
imperfective/perfective-domain contrast in place of the full five-way
indicative inventory. Post-syntactic impoverishment regulates the realisation
of this reduced configuration, as described in
§\,\ref{sec:vm-mood-asp}. Every finite verb, indicative or not, therefore
carries an Asp.Align exponent. The imperative, prohibitive, and hortative
together form the directive domain, a unified subsystem whose case properties
follow from directive orientation within the reduced Asp.Align domain;
it is treated as a group beginning in §\,\ref{sec:mood-directive}.

\subsection{Indicative}

The indicative is the unmarked mood. It alone hosts the full five-way aspect
system and the associated split-ergative alignment described in
Section~\ref{sec:aspect}. Non-indicative moods carry a reduced two-way
aspect-domain contrast in Asp.Align (§\,\ref{sec:vm-mood-asp}), which collapses
the five-way viewpoint inventory; directive Mood additionally fixes the
orientation described below. Overt antipassive morphology is licensed only
with indicative ERG-aspect features, including when that configuration is
embedded under participant nominalisation.

\subsection{The directive domain}\label{sec:mood-directive}

The imperative, the prohibitive, and the hortative form a unified directive
domain: a compact finite subsystem used for commands, negative commands,
suggestions, proposals, soft requests, permissive uses, and wishes. Where the
indicative provides the full five-way aspect system and the split-ergative
alignment described in Section~\ref{sec:aspect}, directive moods carry a reduced
two-way aspect-domain contrast in Asp.Align: the imperfective domain
($\varnothing$ exponent) and the perfective domain
(\ird{-ov-} exponent), both subject to directive orientation. The directive mood suffix follows
reduced Asp.Align, outside any applicative or causative morphology
(§\,\ref{sec:vm-mood-asp}).

Directive Mood selects the theme as the direct-case participant in both
aspectual domains, while an agent or causer with a distinct direct participant
takes transitive case. With no theme, an applied participant or causee can be
selected; an underived intransitive retains its sole argument in direct case.
The understood subject of a command is normally recoverable from the speech
situation, although its omission does not alter this case frame. Applied
participants and causees must be overt in directives; when a theme is retained,
they take indirect case. Antipassive morphology is excluded by a Mood--Voice
licensing restriction (§\,\ref{sec:mood-directive-valency}).

\subsection{Imperative and prohibitive}\label{sec:mood-imp}

The imperative (\ird{-ím}) is the ordinary mood for direct commands,
instructions, warnings, and immediate discourse-management directives addressed
to the hearer. The prohibitive (\ird{-é}) is its dedicated negative counterpart,
used for negative commands and preventative directives. The two are treated
together because they share the same null-subject pattern and the same case
geometry.

The subject of both the imperative and the prohibitive is second person and is
syntactically understood but ordinarily unexpressed; no overt pronoun is
required in fluent discourse. The theme or patient of the clause occupies the
direct case as the clause's single core direct nominal:

\pex
\a
\begingl
\gla Trava piaštím.//
\glb bread.\Dir{} eat-\Imp{}//
\glft \trsl{Eat the bread.}//
\endgl
\a
\begingl
\gla Tak hrabním.//
\glb \Dem{}.\Prox{}.\Dir{} hear-\Imp{}//
\glft \trsl{Hear this.}//
\endgl
\xe

The prohibitive directs the addressee to refrain from the action and excludes
the indicative negative marker \ird{za-}. It is analysed as an exponent of
negative Mood (§\,\ref{sec:vm-moodp}):

\pex
\a
\begingl
\gla Trava piašté.//
\glb bread.\Dir{} eat-\Proh{}//
\glft \trsl{Do not eat the bread.}//
\endgl
\xe

Overt second-person subjects are permitted in both the imperative and the
prohibitive, but they are pragmatically marked, expressing foregrounding,
contrast, insistence, or rebuke toward the addressee. The second-person singular
pronoun is \ird{já} in the direct and \ird{jí} in the transitive; the
second-person plural is \ird{tová} in the direct and \ird{tvy} in the
transitive. In a transitive imperative or prohibitive clause where an overt
direct object is already present, an overt speech-act participant subject
appears in the transitive case rather than the direct, so that the affected
participant may remain the clause's sole direct nominal:

\pex
\a
\begingl
\gla Jí trava piaštím.//
\glb 2\textsc{sg}.\Trs{} bread.\Dir{} eat-\Imp{}//
\glft \trsl{You eat the bread.} \lingcontext{marked; overt 2\textsc{sg} subject in \Trs{}}//
\endgl
\a
\begingl
\gla Tvy trava piašté.//
\glb 2\textsc{pl}.\Trs{} bread.\Dir{} eat-\Proh{}//
\glft \trsl{You lot do not eat the bread.} \lingcontext{marked; overt 2\textsc{pl} subject in \Trs{}}//
\endgl
\xe

\noindent The agent's transitive case belongs to the directive frame with a
separately selected theme; it is retained when the agent is silent. An
underived intransitive has no distinct theme, so its sole participant takes
direct case.

\subsection{Hortative}\label{sec:mood-hort}

The hortative covers a broader range of directive uses than the imperative:
inclusive proposals of the type \textit{let us}, softened direct commands,
permissive uses of the type \textit{let the child come}, and optative
expressions of wish or hope. It serves as the default polite register in
ordinary speech, typically preferred over the imperative in everyday
conversation. The hortative takes two forms: the positive \ird{-ká} and the
negative \ird{-ku}. Like the imperative and prohibitive, it excludes the
ordinary aspect suffixes and operates within the single-direct principle of the
directive domain.

The hortative defaults to a null subject whose interpretation is supplied by
clause type and context. In proposal uses the null subject is most readily read
as first person plural, inclusive of the addressee; in softened directive uses
the same null subject may receive a second-person reading. Null-subject
hortatives are the unmarked form:

\pex
\a
\begingl
\gla Piaštká.//
\glb eat-\Hort{}//
\glft \trsl{Let's eat.} \lingcontext{null subject; inclusive proposal}//
\endgl
\a
\begingl
\gla Tak hrabnká.//
\glb \Dem{}.\Prox{}.\Dir{} hear-\Hort{}//
\glft \trsl{Please hear this.} \lingcontext{null subject; softened directive with direct object}//
\endgl
\xe

The negative hortative \ird{-ku} is used for proposals to refrain, cautious
collective restraint, and mild negative requests, with the suffix expressing
negative hortative Mood and excluding the indicative marker \ird{za-}:

\pex
\a
\begingl
\gla Tak piaštku.//
\glb \Dem{}.\Prox{}.\Dir{} eat-\Neg{}\Hort{}//
\glft \trsl{Let's not eat this.} \lingcontext{negative hortative; collective restraint}//
\endgl
\xe

Overt subjects in the hortative fall into two structurally distinct classes. The
first class comprises first-person plural and second-person subjects in
proposal or softened-command uses. As in the imperative and prohibitive,
these subjects take transitive case when the clause licenses a distinct
participant as direct:

\pex
\a
\begingl
\gla Mně trava piaštká.//
\glb 1\textsc{pl}.\Trs{} bread.\Dir{} eat-\Hort{}//
\glft \trsl{Let us eat the bread.} \lingcontext{overt 1\textsc{pl} subject in \Trs{}}//
\endgl
\a
\begingl
\gla Jí tak hrabnká.//
\glb 2\textsc{sg}.\Trs{} \Dem{}.\Prox{}.\Dir{} hear-\Hort{}//
\glft \trsl{Please, you hear this.} \lingcontext{marked; overt 2\textsc{sg} subject in \Trs{}}//
\endgl
\xe

The second class comprises overt lexical subjects in optative or permissive
hortatives, where the clause expresses a wish or permission directed toward a
third participant rather than a command to the addressee. Such clauses are
typically intransitive. Their lexical subject occupies the direct case as the
clause's single core direct nominal --- the same position that in transitive
directive clauses would be held by the theme or patient. As the sole argument, it takes
direct case:

\pex
\a
\begingl
\gla Byl ščenká.//
\glb child.\Dir{} arrive-\Hort{}//
\glft \trsl{May the child come.} \lingcontext{intransitive optative; lexical subject in \Dir{}}//
\endgl
\a
\begingl
\gla Je ščenká.//
\glb \Dem{}.\Dist{}.\Dir{} arrive-\Hort{}//
\glft \trsl{May that come to pass.} \lingcontext{intransitive optative; distal demonstrative in \Dir{}}//
\endgl
\xe

\noindent A hortative subject thus receives the case licensed by its argument
frame: transitive with a distinct direct-case theme, direct as the sole
participant of an intransitive. The contrast between commands addressed to
speech-act participants and wishes concerning a third participant does not
itself determine case.

\subsection{Case and subject expression in directive clauses}\label{sec:mood-directive-case}

Directive clauses assign direct case to the theme in both aspectual domains.
An agent or causer takes transitive case when a distinct participant is selected
as direct, whereas the sole argument of an underived intransitive remains
direct. The orientation is imposed by directive Mood and therefore persists
whether the subject is pronounced or omitted. The understood second-person
subject of an imperative or prohibitive is normally silent because the speech
situation identifies it; an overt pronoun can contrast or insist on the
addressee without changing the theme's case:

\pex\label{ex:val-directive-simple}
\a
\begingl
\gla Trava piaštím.//
\glb bread.\Dir{} eat-\Imp{}//
\glft \trsl{Eat the bread.} \lingcontext{null subject; theme in \Dir{}}//
\endgl
\a
\begingl
\gla Jí trava piaštím.//
\glb 2\textsc{sg}.\Trs{} bread.\Dir{} eat-\Imp{}//
\glft \trsl{You eat the bread.} \lingcontext{overt subject in \Trs{}; theme in \Dir{}}//
\endgl
\xe

\noindent An overt lexical subject of an intransitive optative or permissive
hortative takes direct case as the clause's sole argument. The transitive
marking of the agent in the second example thus belongs to a clause with a
distinct selected theme; it is independent of grammatical person and of the
choice to pronounce the agent.

\subsection{Valency-changing morphology in directive clauses}\label{sec:mood-directive-valency}

Applicative and causative morphology are compatible with the directive moods.
They precede reduced Asp.Align and the directive mood suffix, so that the
IPFV forms place Mood directly after the derivational suffix, whereas PFV
\ird{-ov-} intervenes before Mood (§\,\ref{sec:vm-mood-asp}).

The applicative (\ird{-éb-}) is licensed in directive clauses only if the
beneficiary is overtly realised. With a retained theme, the beneficiary is an
indirect-case DP, while the theme is direct and any overt agent is transitive.
This Mood-conditioned selection supersedes the simple indicative applicative
frame. A \ird{my} phrase expresses the separate postpositional strategy:

\pex\label{ex:val-directive-appl}
\a
\begingl
\gla Trava bylam piaštébím.//
\glb bread.\Dir{} child.\Ind{} eat-\Appl{}-\Imp{}//
\glft \trsl{Eat the bread for the child's benefit.} \lingcontext{applicative; beneficiary in \Ind{}}//
\endgl
\a
\begingl
\gla Jí trava bylam piaštébím.//
\glb 2\textsc{sg}.\Trs{} bread.\Dir{} child.\Ind{} eat-\Appl{}-\Imp{}//
\glft \trsl{You eat the bread for the child's benefit.} \lingcontext{overt subject in \Trs{}}//
\endgl
\xe

\noindent Without an overt beneficiary the applicative is ungrammatical in
directive clauses:

\pex
\rec{trava piaštébím} \trsl{intended: `Eat bread for someone's benefit.'} \lingcontext{applicative without beneficiary}
\xe

The rhotic causative likewise requires an overtly realised causee in
directive clauses. With a retained lexical theme, the causee takes indirect case
and the theme remains direct; in a causative of an underived intransitive, the
causee itself is direct. An overt causer is transitive in either frame.
Compact \ird{-r-} precedes vowel-initial \ird{-ím} and \ird{-é},
whereas expanded \ird{-er-} precedes consonant-initial hortative
\ird{-ká} and \ird{-ku}: \ird{piaštérká}, \ird{piaštérku}.
When PFV \ird{-ov-} intervenes, the compact grade is retained, as in
\ird{piáštrovím} and \ird{piáštrovká}. Lexically expanded stems,
including statives in \ird{-n-}, retain \ird{-er-}
(§\,\ref{sec:causative-grades}):

\pex\label{ex:val-directive-caus}
\a
\begingl
\gla Trava bylam piáštrím.//
\glb bread.\Dir{} child.\Ind{} eat-\Caus{}-\Imp{}//
\glft \trsl{Make the child eat the bread.} \lingcontext{causative; causee in \Ind{}}//
\endgl
\a
\begingl
\gla Jí trava bylam piáštrím.//
\glb 2\textsc{sg}.\Trs{} bread.\Dir{} child.\Ind{} eat-\Caus{}-\Imp{}//
\glft \trsl{You make the child eat the bread.} \lingcontext{overt subject in \Trs{}}//
\endgl
\xe

With an intransitive base, the overt causee itself takes direct case. The
following constructed pair holds that selection constant across the two
reduced aspectual values; the PFV exponent specifies a bounded directive:

\pex\label{ex:val-directive-caus-intr}
\a
\begingl
\gla Byl zemácrím.//
\glb child.\Dir{} run-\Caus{}-\Imp{}//
\glft \trsl{Make the child run.} \lingcontext{IPFV directive}//
\endgl
\a
\begingl
\gla Byl zemácrovím.//
\glb child.\Dir{} run-\Caus{}-\textsc{pfv.dom}-\Imp{}//
\glft \trsl{Make the child run.} \lingcontext{bounded PFV directive}//
\endgl
\xe

\noindent Without an overt causee the causative is ungrammatical in directive clauses:

\pex
\rec{trava piáštrím} \trsl{intended: `Make someone eat the bread.'} \lingcontext{causative without causee}
\xe

The requirement for overt applied participants and causees is a constructional
restriction of directives. Since several indirect-case DPs may occur together,
a directive causative of an applicative can express the theme in direct case
and both beneficiary and causee in indirect case. When the base has no theme,
the highest eligible participant is direct: the applied participant takes
precedence over a causee. Morphological argument introduction consequently
remains distinct from the selection imposed by higher Mood.

Antipassive morphology is excluded from all directive moods by a
language-specific licensing restriction. Its availability in the indicative
ERG domain does not extend to directive PFV \ird{-ov-}; neither a silent
subject nor a reduced aspect inventory would independently explain this ban.
Commands with an unexpressed patient use the ordinary directive construction,
without overt antipassive morphology.

\subsection{Irrealis}

Affirmative irrealis \ird{-ěv} combines its yod consonantal grade with
local quantitative selection. Where the lexical root is exposed, a
licensed mid-vowel root selects raised-strong G3: \ird{zděk-}
\trsl{bring} gives \ird{zdíčev}. Root /ɛ/ becomes long /iː/,
independently of the velar correspondence /k/ $\rightarrow$ /tʃ/.
Low-vowel \ird{piašt-} instead supplies strong \ird{piáščev},
with quantity on /a/ and a listed cluster alternant. Its consonantal
identity with the antipassive does not insert antipassive Voice into
the irrealis.

\pex
\a \ird{zděk-} + \ird{-ěv} $\rightarrow$ \ird{zdíčev}
\lingcontext{root raised-strong grade; independent yod velar output}
\a \ird{piašt-} + \ird{-ěv} $\rightarrow$ \ird{piáščev}
\lingcontext{root quantity strength; listed modal cluster stem}
\xe

Intervening valency or orientation closes the root domain to this later
Mood selector. With compact Caus, \ird{piáštr-} retains the G1
already supplied by Caus, but irrealis does not assign a new root grade:
\ird{piáštrev} preserves /r/ and takes plain \ird{e}. Expanded
causatives can bear local strength on their own vowel; an already long
\ird{-ér-} retains that quantity, without raising the root through Caus.
The yod consonantal map at a rhotic edge leaves /r/ unchanged.

Reduced perfective Asp supplies its own accessible nucleus. Affirmative
irrealis selects strong \ird{-óv-} from \ird{-ov-}, then applies
PAL to its final /v/, giving \ird{piaštóvěv}. The lexical root
remains short and unmutated. The same local selection gives
\ird{piáštróvěv} after compact Caus and
\ird{piáštranóvěv} where causative-conditioned orientation
\ird{-an-} is present. The retained causative quantity and newly
selected Asp quantity belong to different domains, so their coexistence
does not license rootward reassignment or global shortening.

Negative irrealis \ird{-ir} selects the restricted velar-front
consonantal grade without a new root quantitative grade. Thus
\ird{zděcir} retains the mid vowel, and \ird{piaštir} retains
the short root and plain coronal. It also leaves reduced Asp
\ird{-ov-} short: \ird{piaštovir},
\ird{piáštrovir}, and \ird{piáštranovir}. Previously selected
Caus or lexical quantity remains intact. Negative Mood realises
negation without \ird{za-}; it does not replace the independently
licensed inner stem.

\pex
\a \ird{piaštir} \trsl{would/might not eat; irrealis non-eating}
\a \ird{zděcir} \trsl{would not bring}
\lingcontext{restricted k $\rightarrow$ c; no raised-strong root grade}
\xe

The affirmative irrealis is compatible with the converb-mood
constructions described in Chapter~\ref{ch:complex} and serves as the
main-clause mood in concessive constructions. Its distribution does
not alter the distinction between an exposed root and the edge of an
already assembled complex.


\subsection{Formal analysis: the \synhead{IP} domain}\label{sec:vm-ip}

\subsubsection{MoodP}
\label{sec:vm-moodp}

The mood suffixes follow the aspect-alignment exponent in the verbal complex.
Syntactically, Mood selects Asp directly in affirmative clauses and in negative
non-indicative clauses, while indicative clausal negation introduces \synhead{NegP}
between them (§\,\ref{sec:vm-neg}). In the indicative, \synhead{MoodP}
is present but has a null exponent. Non-indicative Mood licenses the coarser
$[\pm\textsc{pfv-domain}]$ opposition in the syntax, together with its
constructional orientation requirements. Post-syntactic impoverishment
restricts the features available for Vocabulary Insertion at Asp; it neither
assigns case nor changes the interpretation already established by the
syntactic configuration (§\,\ref{sec:vm-mood-asp}). \synhead{AspP} is thus
present throughout the finite paradigm, although the licensed aspectual
contrasts and their exponents vary with Mood.

Mood scopes over the aspectual predicate and conditions the reduced
Asp.Align paradigm. Its affirmative and negative suffixes are competing
Vocabulary Items at the \synhead{MoodP} head:

\begin{center}
\begin{tabular}{l l l}
\toprule
\textsc{mood} & \textsc{affirmative vi} & \textsc{negative vi} \\
\midrule
Indicative   & $\varnothing$      & ---\quad (\synhead{NegP} projected; \ird{za-}) \\
Imperative   & \ird{-ím}        & \ird{-é} \hspace{3pt}(prohibitive) \\
Hortative    & \ird{-ká}        & \ird{-ku} \\
Irrealis     & \ird{-ěv}        & \ird{-ir} \\
\bottomrule
\end{tabular}
\end{center}

\medskip
\noindent Negative non-indicative clauses are analysed with a Mood head bearing
both a mood specification and $[u\textsc{neg}]$, realised by \ird{-é},
\ird{-ku}, or \ird{-ir}. Negative Mood selects Asp directly, whereas
indicative clausal negation contains a separate Neg head below Mood.
Vocabulary Insertion supplies the exponents of these feature bundles after
the syntactic structure has been assembled \cite{embicknoyer2007}. The
absence of a separate Neg node in negative non-indicative clauses accounts for
the exclusion of \ird{za-}; an alternative derivation through post-syntactic
fusion of Neg and Mood would yield the same morphological pattern.

A silent $\mathrm{Op}_{\neg}[i\textsc{neg}]$ in the mood domain licenses
the negative feature of Mood together with any negative indefinites in that
finite clause, extending the concord dependency of
§\,\ref{sec:vm-negative-concord} to \ird{-é}, \ird{-ku}, and \ird{-ir}.
No additional Neg head is required: the overt negative mood suffix satisfies
the obligatory morphological expression of negative polarity. Concord fixes
the number of semantic negations at one within this dependency, without
determining the relative semantic scope of negation and irrealis modality.

The prohibitive \ird{Trava piašté} \trsl{Do not eat the bread} directs the
addressee to refrain from eating, combining directive force with negative
propositional content \cite{han2001}. The negative hortative similarly
proposes or requests restraint, while \ird{piaštir} \trsl{would/might not eat}
expresses the non-occurrence of eating in a hypothetical situation. The
location of the negative feature at Mood is a syntactic analysis of the
suffixal paradigm; the relative semantic scope of negation and irrealis
modality is undocumented.

The negative irrealis \ird{-ir} triggers \ird{k $\rightarrow$ c} in
\ird{k}-final stems, without the \ird{č}-output or vowel raising associated
with affirmative \ird{-ěv}. Thus \ird{zděk-} \trsl{bring} forms
\ird{zděc-ir}, while \ird{piašt-} forms \ird{piaštir}; the forms
\ird{*zděkir} and \ird{*zdíčir} are excluded. The restricted velar-front
grade distinguishes \ird{-ir} from yod-grade \ird{-ěv} and plain-grade
\ird{-ic}.


\subsubsection{Inner evidential (\synhead{EvidP}$_\text{inner}$)}
\label{sec:vm-evid}

An evidential suffix occupies the outermost position in the finite inflectional
spine, immediately inside the complement linker \ird{-e}. This \term{inner
evidential} (\synhead{EvidP}$_\text{inner}$) is distinct from the clausal
evidential and speech-act operators in the left periphery
(§\,\ref{sec:th-particle-framework}), which are Zone~I particles occupying
\synhead{EvidP}$_\text{outer}$ and \synhead{SpeechActP} in the clause spine.

The indirect evidential marks the source of knowledge for the proposition
expressed by the finite verb. Its default exponent is \ird{-al}, with
\ird{-l} selected after a phonologically vowel-final host. The conditioning
depends on the final segment of the fully formed finite host, including any
mood exponent: consonant-final \ird{piaštník} takes \ird{-al}, whereas
vowel-final \ird{piašta} and \ird{piaštíme} take \ird{-l}. A long vowel
or a final diphthongal nucleus also selects \ird{-l}. Selection precedes
coda vocalisation: the /v/ of retrospective \ird{-av} is a consonant
and therefore selects \ird{-al}, even though its free word-final
realisation is [w]. The host retains its vowel quantity; this attachment
requires neither contraction nor an inserted \ird{v}.

The evidential selects the plain grade at its attachment boundary. It adds
no root mutation, vocalic stem grade, or change of orientation and leaves
the finite host's aspect and mood exponents intact. It supplies neither
tense nor person, gender, or number agreement. Both allomorphs belong to
the host's lexical prosodic word and preserve its initial stress.
The order of the outer suffixes is:

\begin{center}
\textsc{stem} + \synhead{Asp.Align} + (\textsc{mood}) +
(\textsc{ind.evid}) + \ird{-e} + \ird{po}
\end{center}

\noindent The corresponding syntactic projection is above
\synhead{MoodP} and constitutes the outermost layer of the finite inflectional
spine. The boundary with the clausal \synhead{EvidP}$_\text{outer}$ (a
left-peripheral projection) must be kept clear: the inner evidential marks the
\emph{source of knowledge} as it bears on the proposition expressed by the
verbal complex; the outer \synhead{EvidP} particles (Zone~I) operate at the
clause-level and take scope over the full proposition including its
discourse-relational frame.

\begin{table}[htbp]
\footnotesize
\caption{Indirect evidential forms across the five indicative aspects.}
\label{tab:inner-evidential-paradigm}
\begin{tabularx}{\textwidth}{L L L L}
\toprule
{\sc aspect} & {\sc finite host} & {\sc indirect evidential} & {\sc linked evidential} \\
\midrule
Perfective & \ird{piaštník} & \ird{piaštníkal} & \ird{piaštníkaľe} \\
Retrospective & \ird{piaštnáv} & \ird{piaštnával} & \ird{piaštnávaľe} \\
Continuous & \ird{piašta} & \ird{piaštal} & \ird{piaštaľe} \\
Progressive & \ird{piaštíme} & \ird{piaštímel} & \ird{piaštímeľe} \\
Contemplative & \ird{piaštáž} & \ird{piaštážal} & \ird{piaštážaľe} \\
\bottomrule
\end{tabularx}
\end{table}

The forms in Table~\ref{tab:inner-evidential-paradigm} retain the ordinary
transitive orientation appropriate to each aspect. The retrospective
contains long \ird{-náv}, so that its evidential ending is
\ird{-nával}. Strong antipassive \ird{piáščav} instead preserves
short \ird{-av} in \ird{piáščaval}. The \ird{a} of continuous \ird{piašta} belongs to
Asp.Align: \ird{piašta-l} yields \ird{piaštal}, without deleting an
aspect vowel or adding a second one.

\paragraph{The outer linker.}
When \ird{-e} follows the evidential, its front-grade mapping targets
evidential-final /l/, which selects [ʎ], written \ird{ľ}. Thus
\ird{-al-e} gives \ird{-aľe}, and postvocalic \ird{-l-e} gives
\ird{-ľe}. The preceding finite host remains intact in both:
\ird{piaštíme-l-e} yields \ird{piaštímeľe}. The \ird{e} before
\ird{ľ} in the sequence \ird{-eľe} is the progressive's own final vowel.
Syllabification follows the complete suffixal sequence, so a free
evidential's final hard /l/ vocalises to [w], whereas the linked soft
lateral remains [ʎ] (§\,\ref{sec:l-vocalisation}).

This outer boundary does not reach through the evidential to select
inner linker allomorphs. Perfective \ird{piášček} accordingly forms
\ird{piáščekal} and linked \ird{piáščekaľe}, preserving \ird{-ek};
its directly linked counterpart \ird{piáščice} belongs to the distinct
construction without intervening Evid. Likewise, \ird{piaštímeľe}
retains Class~B \ird{-íme}, whereas direct linking without Evid selects
\ird{piaštímňe}. The contrast follows the local boundary of each
selecting exponent (§\,\ref{sec:boundary-mapping}).

\paragraph{Combinations with voice, causative, and mood.}
The suffix attaches outside the independently licensed finite construction.
Table~\ref{tab:evidential-combinations} gives constructed combinations;
the causative forms preserve both the rhotic exponent and outer
orientation \ird{-an-} where required. Mood supplies the input to
evidential allomorphy, so prohibitive \ird{piašté} selects \ird{-l}
while imperative \ird{piaštím} selects \ird{-al}. Their long vowels
remain long. The directive rows illustrate exponence where a combination
of Evid and directive Mood is independently licensed; they do not establish
its semantic or pragmatic distribution. The choice of \ird{-al} or
\ird{-l} alone does not settle whether such a combination is available.

\begin{table}[htbp]
\footnotesize
\caption{Constructed combinations of the evidential with finite morphology.}
\label{tab:evidential-combinations}
\begin{tabularx}{\textwidth}{L L L L}
\toprule
{\sc construction} & {\sc finite host} & {\sc indirect evidential} & {\sc linked evidential} \\
\midrule
Antipassive perfective & \ird{piášček} & \ird{piáščekal} & \ird{piáščekaľe} \\
Antipassive retrospective & \ird{piáščav} & \ird{piáščaval} & \ird{piáščavaľe} \\
Passive progressive & \ird{piaštníme} & \ird{piaštnímel} & \ird{piaštnímeľe} \\
Causative progressive & \ird{piáštrime} & \ird{piáštrimel} & \ird{piáštrimeľe} \\
Causative perfective & \ird{piáštranik} & \ird{piáštranikal} & \ird{piáštranikaľe} \\
Irrealis & \ird{piáščev} & \ird{piáščeval} & \ird{piáščevaľe} \\
Negative irrealis & \ird{piaštir} & \ird{piaštiral} & \ird{piaštiraľe} \\
Imperative & \ird{piaštím} & \ird{piaštímal} & \ird{piaštímaľe} \\
Prohibitive & \ird{piašté} & \ird{piaštél} & \ird{piaštéľe} \\
\bottomrule
\end{tabularx}
\end{table}

Indicative negation likewise preserves its ordinary form outside the lexical
stem: \ird{zapiaštník} gives \ird{zapiaštníkal} and linked
\ird{zapiaštníkaľe}. Evidential suffixation neither replaces negative
Mood \ird{-ir} with \ird{za-} nor changes the argument frame licensed
by the inner verb.




\section{Non-finite forms and nominalisation}\label{sec:converb-morphology}


Converb suffixes replace the aspect-alignment slot entirely. A verb carrying a
converb suffix accordingly bears no finite viewpoint or mood exponent and
no aspect-conditioned alignment contrast. Its clause nevertheless licenses
nominal case through the converb construction (§\,\ref{sec:vm-converb-case}). The
morphological structure of a converb is:

\medskip
\begin{center}
(\textsc{lex-prefix}) + \Stem{} + (\Appl{}) + (\Caus{}) + \textsc{converb-suffix}
\end{center}
\medskip

The converb suffixes and the stem mutations they occasion are set out in the following table:

\begin{table}[ht]
	\footnotesize
	\caption{Converb suffixes and their phonological effects.}\label{tab:converb-morph}
	\begin{tabularx}{\textwidth}{L L L L}
		\toprule
		{\sc converb} & {\sc suffix(es)} & {\sc stem mutation} & {\sc notes} \\
		\midrule
		Sequential    & \ird{-u}, \ird{-ení} & none & \ird{-ní} after applicative; full \ird{-ení} after causative \\
		Simultaneous/manner & \ird{-ěc} & palatalisation & ordinary consonant mutation; independent stem-grade selection \\
		Temporal (6)  & \ird{-ezak}, \ird{-eškady}, \ird{-eeby}, \ird{-ešhoume}, \ird{-eebymy}, \ird{-esny} & none & all add \ird{-e-} linker \\
		Concessive    & \ird{-ouc} & none & \\
		Conditional   & \ird{-ic}, \ird{-emy}, \ird{-ezmy}, \ird{-ena} & none for \ird{-ic}; see note & \ird{-ic} does not trigger palatalisation \\
		Purposive     & \ird{-it} & irregular; mainly Class I & \\
		Contrastive   & \ird{-ema} & none & \\
		\bottomrule
	\end{tabularx}
\end{table}

The simultaneous/manner converb \ird{-ěc} selects the front grade,
irrealis \ird{-ěv} the yod grade, and antipassive \ird{-aš-} its separate
replacement classes. Conditional \ird{-ic} selects the plain grade
(§\,\ref{sec:cond-nopala} of Chapter~\ref{ch:complex}); its initial vowel
leaves alveolar sibilants unchanged. Sequential \ird{-ení} is also
plain-grade, unlike the homophonous front-coded nominal ending. The shortened
\ird{-ní} is selected immediately after applicative \ird{-éb-}; a following
rhotic causative instead selects full \ird{-ení}, as in \ird{piáštrení}
and \ird{piaštébrení}, beside the alternative in \ird{-u}. Simultaneous
\ird{-ěc} leaves causative /r/ unchanged and is written with plain
\ird{e}: \ird{piáštrec}.
Grade membership must therefore be recorded with each construction.


\subsection{Formal analysis: the syntactic status of converbs}
\label{sec:vm-converbs}

\subsubsection{Converbs as aspect-deficient \synhead{vP}/\synhead{VoiceP} projections}
\label{sec:vm-converb-aspless}

Converbs lack finite Asp.Align and its viewpoint distinctions, while retaining
the lower predicate and any licensed applicative or causative structure.
The analysis treats the converb suffix as the exponent of a converb head,
Conv, with a defective aspectual specification. Its features specify the
linking relation without supplying
the finite IPFV/ERG-domain opposition. Conv selects the lower verbal
predicate directly, excluding the finite perfective/retrospective configuration
that licenses ergative alignment.
The independent absence of outer patient-selecting Voice, including both
ordinary \ird{-n-} and causative-conditioned \ird{-an-}, from the
established converb forms is described in §\,\ref{sec:converb-chains};
its exclusion does not follow from the absence of ergative case alone, since
the same Voice head also occurs in finite nominative-domain passives.

\subsubsection{Case alignment within converb clauses}
\label{sec:vm-converb-case}

Converb clauses license the NOM-domain case frames without finite viewpoint
aspect. An ordinary active selects its agent or sole participant as direct and
assigns a retained patient indirect case; a simple applicative selects the
applied participant, with actor and retained theme indirect, while a higher
causative selects its causer. No converb licenses transitive-case agents.
These are constructional licensing conditions: absence of perfective or
retrospective Asp excludes the ERG frame, but does not by itself determine
which participant is privileged.

\subsubsection{Converb clauses as adjuncts and embedded topics}
\label{sec:vm-converb-adjunct}

In the clause-combining hierarchy, converb clauses function as adverbial
adjuncts to the matrix clause. They are merged as adjuncts to the matrix
\synhead{IP} or \synhead{vP}, not as arguments in Spec or Complement positions.
This explains their preverbal (and pre-matrix) position: as adjuncts to
\synhead{IP}, they precede all material within \synhead{IP} and surface before
the matrix arguments in the linear string, consistent with the head-final
adjunction pattern.

Reference across converb chains follows a default same-pivot relation
 (§\,\ref{sec:switch-reference}). Distinct overt nominals or discourse
 context can support different-pivot readings, without a switch-reference
 clitic. Since an aboutness topic need not be the absolutive pivot, topic
 continuity and this dependency are separate relations. These constructions
 do not by themselves establish the availability of embedded hanging topics
 or high topicalisation in every non-finite clause.

\subsection{The nominalised verb forms}\label{sec:nominalised}


The suffix \ird{-ou} occurs in lexical nouns, event nominals, and
aspect-bearing participant descriptions. Lexical \ird{piaštou} \trsl{food}
denotes a substance conventionally associated with eating, whereas eventive
\ird{piaštou} \trsl{eating} denotes an eventuality and admits event-related
dependents, as in \ird{Marcí travěte piaštou} \trsl{Marek's eating of the
bread}. Participant \ird{piáščikou} \trsl{the one who ate} denotes an
individual characterised by participation in an event. The shared exponent
therefore does not entail a single internal structure. Event, result, and
participant readings commonly coexist in nominalisation systems
\cite{lieber2018}; the structures distinguished below constitute the analysis
of Iridian adopted here.

Lexical nominalisation combines a root or a small verbal constituent with
$n_{\mathrm{lex}}$, whose interpretation may be stored with the derivative.
Event nominalisation combines an event predicate with $n_{\mathrm{evt}}$,
without selecting an individual participant as its referent. Its established
forms do not require the five-way indicative aspect distinction. Participant
$n_{\mathrm{part}}$ selects \synhead{AspP}, containing
the inner valency domain, optional \synhead{CausP}, and the outer
orientation head.
Embedded Asp retains viewpoint and participant-selection features, although
finite case valuation is inactive: the nominal domain licenses the indirect
case on its dependents (§\,\ref{sec:nm-derived-gerund}). Lexical agent nouns
in \ird{-(h)ár} and \ird{-ist} (Chapter~\ref{ch:derivational}) classify a
person by a characteristic activity or occupation; aspect-bearing \ird{-ou}
forms describe a participant in an aspectually construed situation. This
separation accords with Baker and Vinokurova's distinction between lexical
agent nominalisation and the richer verbal structure of reduced relatives
\cite{bakervinokurova2009}.

\paragraph{Aspect and participant selection.}
Participant nominalisation targets the participant privileged by the complete
voice--aspect configuration, rather than assigning an invariant orientation to
a lexical transitive stem. With a transitive base in the perfective or
retrospective, antipassive Voice privileges the agent, while \ird{-n-}
privileges the patient. In the continuous, progressive, and contemplative,
active Voice already privileges the agent, and patient-oriented \ird{-n-}
yields the passive configuration (§\,\ref{sec:vm-n}). Antipassive morphology
remains excluded from these three aspects, including under nominalisation.

\begin{table}[htbp]
\footnotesize
\caption{Participant nominals of \ird{piašt-} \trsl{eat}. English glosses
in the text give customary contextual readings.}\label{tab:participant-nominals}
\begin{tabularx}{\textwidth}{l L L}
\toprule
Aspect & Agent-oriented & Patient-oriented \\
\midrule
Perfective & \ird{piáščikou} & \ird{piaštníkou} \\
Retrospective & \ird{piáščavou} & \ird{piaštnávou} \\
Continuous & \ird{piaštova} & \ird{piaštnova} \\
Progressive & \ird{piaštímou} & \ird{piaštnímou} \\
Contemplative & \ird{piaštážou} & \ird{piaštnážou} \\
\bottomrule
\end{tabularx}
\end{table}

The perfective selects quantity-preserving bound
\ird{-ik-}/\ird{-ík-} before nominalising \ird{-ou}.
Thus short \ird{piášček} corresponds to \ird{piáščikou},
whereas suffix-strong \ird{piaštník} gives \ird{piaštníkou}.
Retrospective \ird{-av}/\ird{-áv} similarly gives
\ird{-avou}/\ird{-ávou}. Continuous Asp together with
$n_{\mathrm{part}}$ has joint direct-singular \ird{-ova}.
Progressive \ird{-ime}/\ird{-íme} corresponds to
\ird{-imou}/\ird{-ímou}, preserving the selected quantity.
Contemplative \ird{-áž}/\ird{-až} combines with \ird{-ou}
as \ird{-ážou}/\ird{-ažou}. These Vocabulary Insertion conditions specify constructional
allomorphy; neither \ird{-ova} nor \ird{-imou} follows from a general
vowel-deletion or root-raising rule. Acute accents retain their ordinary value
as quantity marks rather than stress marks (§\,\ref{sec:stress}).

Perfective \ird{piáščikou} describes the agent of a bounded eating event,
and retrospective \ird{piáščavou} an agent whose participation precedes the
relevant reference point. Their customary translations \trsl{the one who ate}
and \trsl{the one who has eaten} do not add grammatical tense. Continuous
\ird{piaštova} permits habitual or ongoing interpretation, progressive
\ird{piaštímou} profiles participation during the event's development, and
contemplative \ird{piaštážou} describes a prospective or intended
participant. The latter may denote an existing individual although the eating
has not occurred.

For the perfective agent form, a schematic interpretation is
\[
 \lambda x\,\exists e\,[\operatorname{eat}(e)\land
 \operatorname{Agent}(e,x)\land\operatorname{PF}(e,R)],
\]
where $R$ is the contextual reference frame. Eventive \ird{piaštou}, by
comparison, supplies a property of eating events rather than a property of
eaters. The existential closure illustrated for the perfective does not extend
to an assertion of an actual event with the contemplative. Neither
$n_{\mathrm{part}}$ nor \ird{-ou} encodes uniqueness or definiteness;
anchoring and the construction in which the nominal occurs determine its
referential interpretation.

Applicative nominalisation selects the applied participant, as in
\ird{piaštébikou} \trsl{the one for whose benefit eating took place}.
Causativisation establishes a higher orientation domain, so that participant
nominalisation of an active causative selects the causer in NOM aspects and
the retained theme, otherwise the applied participant, otherwise the causee in
ERG aspects. A NOM-domain causative bearing passive \ird{-an-} selects the
same embedded participant hierarchy instead of the causer
(§\,\ref{sec:vm-n}). The suffixal sequence is root--inner valency--causative--orientation--Asp--nominalisation;
the nominaliser adds no further promotion operation. Thus an inner applicative
need not supply the referent of a causative participant nominal.

\paragraph{The nominal case domain.}
A non-privileged participant may be expressed by an indirect-case nominal
dependent, whose relation to the embedded event is licensed in the nominal
domain. In the analysis adopted here, a silent nominal relational head
$\mathrm{R}_{\theta}$ introduces the dependent and adds an event relation
such as Agent or Theme to the participant description before event closure.
Its projection adjoins to the event predicate within the nominaliser's
complement, preserving the selected \synhead{AspP} category in participant
nominals. The relational head licenses an indirect \synhead{KP}; the
dependent's K head realises case, as in other nominal dependency relations
(§\,\ref{sec:nm-poss-struct}). This relation neither assigns a finite verbal
case nor projects another core argument inside \synhead{VP}. Thus a patient
dependent of an antipassive-based nominal
restricts the eating event by Theme$(e,y)$, although antipassive
\synhead{VoiceP} itself contains no overt patient argument. This constructional
relation draws on the lexical event meaning; it does not reverse
antipassivisation. Finite \ird{*Marek travě piášček} consequently remains
ungrammatical (§\,\ref{sec:vm-antipassive-desc}). The restriction is specific
to Iridian finite antipassives: cross-linguistically, the availability of an
oblique patient varies among antipassive constructions \cite{heaton2020}.

A thematic dependent differs from an ordinary possessor, which relates an
individual to the nominal referent without necessarily identifying an event
role, and from an adjunct such as an instrumental postpositional phrase,
whose relation is contributed by the postposition. All may contain indirect
case, so that case alone identifies neither their interpretation nor their
position in the nominal structure.

\paragraph{Negation and the upper boundary.}
The aspect-bearing participant nominal and the corresponding \ird{-ový}
relative select \synhead{AspP}, below clausal \synhead{NegP}. Internal verbal
\ird{za-} is therefore excluded from these constructions; it cannot acquire
scope inside them merely by attaching to their surface stem. Negation with
\ird{staly} rejects the matrix predication, while \ird{zad-} denies the
existence of an individual satisfying the positive participant description.
Neither expresses \trsl{the one who did not eat}. A finite negative clause
provides the established way to assert non-participation; no internally
negative \ird{-ou} or \ird{-ový} paradigm is licensed here. The smaller
lexical and event nominalisations likewise supply no productive clausal
negation domain. Strict negative concord remains a property of the appropriate
clausal domain, rather than a mechanism for extending nominalisation above
\synhead{NegP}.

\subsubsection{Predicative use: participant nominals in identificational equatives}
\label{sec:nmlz-predicative}

Participant nominals may serve as the predicative complement of a zero-copula
equative clause, with the topic \synhead{DP} identified as the entity satisfying
the nominalised description. This construction has the schematic form:

\begin{center}
[\textit{topic-DP}] \ \ [\textit{(nominal dependent.\Ind{})} \ \ \textit{participant-nominal}]
\end{center}

\noindent Compare the base transitive clause with its patient-nominal equative counterpart:

\pex
\a
\begingl
\gla Lobek Janku piaštník.//
\glb apple.\Dir{} Janek.\Trs{} eat-\textsc{trs}-\Pf{}//
\glft \trsl{Janek ate the apple.} \lingcontext{standard ergative clause}//
\endgl
\a
\begingl
\gla Lobek Jancí piaštníkou.//
\glb apple.\Dir{} Janek.\Ind{} eat-\textsc{trs}-\Pf{}-\textsc{nmlz}//
\glft \trsl{The apple is what Janek ate.} \lingcontext{patient nominal as equative predicate}//
\endgl
\xe

\noindent And with an agent-nominal equative:

\ex
\begingl
\gla Janek lúbcí piáščikou.//
\glb Janek.\Dir{} apple.\Ind{} eat-\Antip{}-\Pf{}-\textsc{nmlz}//
\glft \trsl{Janek is the one who ate the apple.} \lingcontext{agent nominal as equative predicate}//
\endgl
\xe

The finite transitive agent \ird{Janku} bears the transitive case,
whereas nominal \ird{Jancí} receives indirect case within the patient
nominal. In the agent nominal, indirect \ird{lúbcí} restricts the event
through the nominal Theme relation just defined; it is not an internal
argument retained by antipassive \synhead{VP}. These dependents receive
nominal case independently of the external topic and independently of whether
the participant nominal functions as an equative predicate or an existential
pivot (Chapter~\ref{ch:simple}).

On the identificational reading illustrated here, the participant nominal
supplies a description against which the topic is identified, with the context
providing the relevant event and participant. This differs from the property
ascription in \ird{Marek doktor} \trsl{Marek is a doctor} and from the finite
assertion \ird{Lobek Janku piaštník} \trsl{Janek ate the apple}, without
requiring that participant nominalisation itself encode unique reference.
Negation of identificational equatives
with \ird{staly} and its semantic contrast with the verbal \ird{za-} are treated
in §\,\ref{sec:ident-equative} and §\,\ref{sec:clausal-negation} of
Chapter~\ref{ch:simple}.


\section{The conjunctive / quotative linker}\label{sec:conjunctive-linker}

Iridian has a verbal linking form in \ird{-e}, used in clause linking and
quotative constructions. The suffix selects the front grade at the
right edge of the fully inflected verb. Its morphological grade is independent
of the fact that its pronounced vowel is /ɛ/.

\pex
\a \ird{piášček} $\rightarrow$ \ird{piáščice}
\a \ird{piáščav} $\rightarrow$ \ird{piáščavě}
\a \ird{mrocna} $\rightarrow$ \ird{mrocne}
\a \ird{piaštíme} $\rightarrow$ \ird{piaštímňe}
\a \ird{mrocnáž} $\rightarrow$ \ird{mrocnáže}
\xe

\begin{table}[ht]
  \footnotesize
  \caption{Representative linker outputs across finite paradigms.}
  \label{tab:linker-paradigm}
  \begin{tabularx}{\textwidth}{l l l l}
    \toprule
    \textsc{mood} & \textsc{domain / category} & \textsc{finite form} & \textsc{linked form} \\
    \midrule
    Indicative & perfective & \ird{piaštník} & \ird{piaštníce} \\
    Indicative & retrospective & \ird{piaštnáv} & \ird{piaštnávě} \\
    Indicative & continuous & \ird{mrocna} & \ird{mrocne} \\
    Indicative & progressive & \ird{piaštíme} & \ird{piaštímňe} \\
    Indicative & contemplative & \ird{mrocnáž} & \ird{mrocnáže} \\
    Imperative & IPFV & \ird{piaštím} & \ird{piaštímě} \\
    Imperative & PFV & \ird{piaštovím} & \ird{piaštovímě} \\
    Prohibitive & IPFV & \ird{piašté} & \ird{piašťe} \\
    Prohibitive & PFV & \ird{piaštové} & \ird{piaštově} \\
    Hortative & IPFV & \ird{piaštká} & \ird{piaštce} \\
    Hortative & PFV & \ird{piaštovká} & \ird{piaštovce} \\
    Negative hortative & IPFV & \ird{piaštku} & \ird{piaštcuje} \\
    Negative hortative & PFV & \ird{piaštovku} & \ird{piaštovcuje} \\
    Irrealis & IPFV & \ird{piáščev} & \ird{piáščevě} \\
    Irrealis & PFV & \ird{piaštóvěv} & \ird{piaštóvěvě} \\
    Negative irrealis & IPFV & \ird{piaštir} & \ird{piaštire} \\
    Negative irrealis & PFV & \ird{piaštovir} & \ird{piaštovire} \\
    \bottomrule
  \end{tabularx}
\end{table}

Consonant-final forms expose the grade directly: /k/ selects /ts/,
input /ts/ selects /tʃ/, and final labials acquire secondary palatalisation.
Explicit \ird{ť\,ď\,ň\,ľ} represent the corresponding soft coronals.
The paradigm also contains listed ending allomorphs, which cannot all be
derived by mutation of a final consonant: \ird{-na/-ne},
\ird{-ime/-imňe}, \ird{-é/-ťe},
\ird{-ká/-ce}, and \ird{-ku/-cuje}. Their extra consonants, vowel
replacements, or quantity relations belong to the individual construction.
Thus linked \ird{mrocne} preserves /n/ within the ending \ird{-ne};
it differs from adverbial \ird{mrocňe}, where PAL \ird{-ě} selects /ɲ/.
Contemplative \ird{mrocnáž} instead gives \ird{mrocnáže} by regular
linker suffixation: PAL preserves the input postalveolar /ʒ/.
The difference between free final [ʃ] and prevocalic [ʒ] follows
ordinary final devoicing (§\,\ref{sec:lenition}).


\subsection{Formal analysis: the conjunctive linker \ird{-e}}
\label{sec:vm-tension-linker}

The linker \ird{-e} appears at the right edge of embedded finite clauses before
\ird{po}/\ird{poby} and also in the conjunctive clause-chaining function
described above. The analysis adopted here assigns these uses a common
clause-edge function, supported by their shared suffix position and
consonant-mutation class. The two environments remain distinguishable
in the syntactic contexts that license the suffix.

\subsubsection{Phonological characterisation}

\paragraph{Trigger class and palatalisation.}
The linker selects the front grade described in
§\,\ref{sec:suffix-classes}. It shares the coronal and labial PAL mappings
with front-coded nominal endings, while its class-assigned velar output /ts/
differs from the /tʃ/ selected by yod-grade irrealis \ird{-ěv}.
Plain-grade conditional \ird{-ic} and restricted velar-front \ird{-ir}
have different specifications. These four constructional classes combine two
instructions: whether the coronal and labial PAL mapping applies, and which
velar output is selected. The linker combines PAL with /k/ $\rightarrow$ /ts/;
irrealis \ird{-ěv} combines PAL with /k/ $\rightarrow$ /tʃ/;
\ird{-ir} selects /k/ $\rightarrow$ /ts/ without PAL; and \ird{-ic}
selects neither mutation. Surface vowel quality does not determine either
instruction. At each licensed boundary the mapping takes the selected input
consonant once; it does not feed its own output back into the same operation.

\paragraph{Segment-specific outputs.}
The grade targets the final consonant of the inflected input at its own
boundary. It does not reach through an intervening aspect or mood exponent
to mutate the thematic consonant of the bare root. Where the evidential
intervenes, that input ends in its /l/, so the linker selects \ird{ľ}
and the inner aspect or mood ending remains intact
(§\,\ref{sec:vm-evid}).

\begin{itemize}
  \item For input /k/, the front grade selects /ts/, spelled \ird{c}.
    In \ird{piášček} $\rightarrow$ \ird{piáščice}, perfective
    \ird{-ik-} is selected first, and its final velar then mutates.
    The resulting /ts/ does not re-enter PAL at the same boundary.
  \item Input /t\,d\,n\,l/ selects [c\,ɟ\,ɲ\,ʎ], written
    \ird{ť\,ď\,ň\,ľ}; input /s\,z\,ts/ selects /ʃ\,ʒ\,tʃ/, written
    \ird{š\,ž\,č}.
    Postalveolars already present within the form are unaffected.
  \item Labials /p\,b\,m\,v/ acquire secondary palatalisation, marked by
    \ird{ě} before /ɛ/, as in \ird{piaštímě} and \ird{piáščevě}.
    Retrospective \ird{piáščav} accordingly gives \ird{piáščavě},
    with prevocalic [ʋʲ] at the linker boundary.
    Rhotic /r/ remains a rhotic: \ird{piaštir} gives \ird{piaštire},
    without obligatory secondary palatalisation. Optional phonetic [rʲ]
    is not an additional PAL output (§\,\ref{sec:soft-coronals}).
\end{itemize}

Consonant mutation precedes coda vocalisation and prevocalic approximant
realisation. The linked soft lateral remains [ʎ], whereas prevocalic /v/ carries its secondary
palatal articulation on [ʋʲ]. Coda /l/ and /v/ vocalisation is determined
by the syllabification of the complete linked form, before final devoicing
applies to its remaining obstruents
(§§\,\ref{sec:l-vocalisation} and~\ref{sec:v-realisation}).

The mappings apply to the boundary input once. Listed ending allomorphs in
Table~\ref{tab:linker-paradigm} supplement these productive correspondences;
they do not extend the rule to every segment found inside an ending.

\paragraph{Selection of the aspect grade.}
In directly linked perfectives without intervening Evid, the linker selects
bound \ird{-ik-}, after which \ird{k} becomes
\ird{c}: \ird{piášček} $\rightarrow$ \ird{piáščik-e}
$\rightarrow$ \ird{piáščice}. The lexical root is unchanged. Transitive
\ird{piaštník} already has the high-vowel exponent and gives
\ird{piaštníce}. This suffixal correspondence differs from root raising in
\ird{zděk-}/\ird{zdíčev}. The perfective suffixal height relation
preserves the class-selected quantity, including long \ird{-ék}/\ird{-ík-};
irrealis G3 changes both height and quantity at its exposed root target.
Neither operation affects every \ird{e/ě} before a front-vocalic suffix.

Other linked aspect and mood forms in Table~\ref{tab:linker-paradigm} contain
construction-specific ending allomorphs. Final-vowel replacement or quantity
change in such pairs as \ird{mrocna}/\ird{mrocne} and
\ird{piašté}/\ird{piašťe} belongs to those listed exponents, not to a
productive shortening or raising of the root.

\paragraph{Listed ending allomorphs.}
Conjunctive forms whose ending correspondence exceeds ordinary mutation are
listed with their aspect or mood paradigm. Thus the \ird{-e} of
\ird{piašťe} corresponds to free prohibitive \ird{-é}, while
\ird{piaštímňe} contains long \ird{-ímňe}, corresponding to
Class~B \ird{-íme}; root-strong or lexical-long classes retain the
short \ird{-ime}/\ird{-imňe} pair. These
forms require construction-specific allomorphs, even though they occupy the same
linker position as productive \ird{-ice}. Their selection requires the
corresponding finite ending at the linker boundary: \ird{piáščekal}
instead yields \ird{piáščekaľe}, and \ird{piaštímel} yields
\ird{piaštímeľe}. The evidential intervenes between Asp and the outer
linker, which therefore does not select inner \ird{-ik-} or
\ird{-imňe}.

\paragraph{Prosodic status.}
The linker \ird{-e} belongs to the same minimal lexical prosodic word as
its inflected host. Once the selected stem grade and the resulting syllables
have been established, primary stress falls on the first syllable of that
domain (§\,\ref{sec:stress}). Outer suffixation consequently does not move
stress to the linker. An integrated prefix remains inside the domain and bears
initial stress where present; an external negative
\ird{za-} does not enter it. Stress neutrality under suffixation therefore
does not imply stress fixed to the bare lexical root.

\subsubsection{Morphosyntactic distribution}

\paragraph{Template position.}
Within the verbal complex, \ird{-e} attaches outside all finite inflectional
morphology. The full template is:

\begin{center}
\Stem{} + \synhead{Asp.Align}.\textsc{align} + (\textsc{mood}) +
(\textsc{evid}) + \ird{-e}
\end{center}

\noindent The linker is therefore exterior to \synhead{EvidP}$_\text{inner}$,
the outermost projection in the finite inflectional spine
(§\,\ref{sec:vm-evid}). This position is significant: if \ird{-e} were an
agreement or mood morpheme, it would be expected inside the inflectional spine.
Its consistent extra-clausal position supports an analysis in which it marks the
edge of the finite complex as it interfaces with a higher projection --- a
\synhead{CP}-boundary marker, not an inflectional head.

\paragraph{Distribution.}
The linker appears in exactly three environments:

\begin{enumerate}
  \item \textit{The finite \ird{po}-complement} (§\,\ref{sec:vm-po-comp}): the
    linker marks the right edge of the embedded finite \synhead{TP} before the
    complement C head \ird{po} or \ird{poby}. It is obligatory here when the
    embedded predicate is verbal; an embedded finite verbal clause without the
    linker is ungrammatical before \ird{po}/\ird{poby}. With zero-copula
    predicates, where no finite verb is present to host \ird{-e}, the support
    form \ird{ce} appears before the complementizer instead. The two
    complementizers attach not to a bare proposition versus a bare
    \textit{whether}-clause, but to assertive versus non-assertive finite
    complements respectively; the non-assertive type includes embedded
    wh-content clauses as well as embedded polar issues (see §\,\ref{sec:poby}
    of Chapter~\ref{ch:complex}).

  \item \textit{The conjunctive chain} (§\,\ref{sec:conjunctive-linker}): the
    linker marks non-final finite clauses in a clause chain where the chained
    clauses are interpreted as coordinated or sequentially related propositions.
    This use resembles the clause-linking \textit{-ko} of Korean or the
    converbal linker of other verb-final languages, but unlike those, the
    Iridian form attaches to a fully finite, fully inflected verb.

  \item \textit{The matrixless \ird{po}-construction}
    (§\,\ref{sec:complementation}): even when the matrix predicate is absent and
    \ird{po} functions as a discourse-level attribution marker, the linker
    \ird{-e} is retained on the finite verb. This shows that \ird{-e} is not
    licensed by the presence of a lexical matrix predicate but by the presence
    of the C-domain element \ird{po} itself.
\end{enumerate}

\paragraph{Non-environments.}
The linker is absent in two structurally revealing contexts. First, independent
finite clauses lack \ird{-e}: when no C head embeds the \synhead{TP}, the linker
is not licensed. Second, converb forms --- the non-finite clause-linking
morphology of simultaneous, sequential, conditional, and other converbs --- do
not take \ird{-e}. This follows naturally: converbs are already non-finite (they
lack Asp.Align and Mood morphology), so they do not participate in the domain of
the finite-clause boundary that \ird{-e} marks.

\subsubsection{Clause-edge analysis}

The grammar adopts a unified \synhead{CP}-edge analysis.
Its core claim is that \ird{-e} spells out the feature \textsc{[+fin, +sub]} ---
roughly, `finite clause under subordination' --- on the embedded \synhead{TP}'s
outermost layer, and that this feature is valued whenever a C head (overt or
null) selects the finite \synhead{TP}. The front grade selected by
\ird{-e} is then part of the PF realisation of this morphosyntactic feature,
parallel to the corresponding nominal mutation class. The exponent's vowel
does not itself derive this specification.

The position of \ird{-e} outside \synhead{EvidP}$_\text{inner}$ in the template
supports this analysis directly: the linker is the last element added to the
verbal complex before the clause exits the finite inflectional domain, which is
exactly where a \synhead{TP}-edge or C-interface marker would be expected.

\paragraph{Locality with the evidential.}
The evidential supplies a consonantal boundary between the inner finite
ending and the linker. The linked sequences \ird{-aľe} and
postvocalic \ird{-ľe} follow from PAL applied to evidential /l/,
without mutation or linker-conditioned allomorphy across Evid
(§\,\ref{sec:vm-evid}). This local realisation is compatible with the clause-edge analysis.

\noindent The shared mutation class and ending selection associate
\ird{-e} with the morphological system of the verb, even though the linker
occupies the clause edge outside finite inflection. These properties support its
analysis as a suffix.


\section{Stative verbs}\label{sec:stative}

A class of stative verbs describes states rather than events. These verbs are
the language's primary means of expressing properties predicatively: where
English often uses an adjective, Iridian normally uses a stative verb. Stative
verbs are true verbs and accordingly require no copula to function as
predicates.

\subsection{Property roots and the verbalising suffix \ird{-n-}}

A small closed class of property roots each denotes an abstract property as a
nominal. The verbalising suffix \ird{-n-} is added to the root to yield a
stative verb stem, which then conjugates as a regular intransitive verb:

\begin{table}[ht]
	\footnotesize
	\caption{Property roots and derived stative verbs.}\label{tab:statroots}
	\begin{tabularx}{\textwidth}{l L l L}
		\toprule
		{\sc root} & {\sc gloss} & {\sc stative stem} & {\sc gloss} \\
		\midrule
		\ird{hrad} & goodness  & \ird{hradn-}  & to be good \\
		\ird{mroc} & beauty    & \ird{mrocn-}  & to be beautiful \\
		\ird{trpa} & badness   & \ird{trpan-}  & to be bad \\
		\ird{svor} & bigness   & \ird{svorn-}  & to be big \\
		\bottomrule
	\end{tabularx}
\end{table}

The derived stative verb conjugates across the full five-way aspect system. The
continuous is closest to an equative property predication; the other aspects
carry characteristically dynamic or inchoative readings. The full paradigm of
\ird{mrocn-} \trsl{be beautiful} illustrates the system:

\pex
\a
\begingl
\gla Nužo mrocna.//
\glb city.\Dir{}.\Def{} be.beautiful-\Cont{}//
\glft \trsl{The city is beautiful.} \lingcontext{continuous: current state}//
\endgl
\a
\begingl
\gla Nužo mrócnime.//
\glb city.\Dir{}.\Def{} be.beautiful-\Prog{}//
\glft \trsl{The city is becoming beautiful.} \lingcontext{progressive: entering the state}//
\endgl
\a
\begingl
\gla Nužo mrócnik.//
\glb city.\Dir{}.\Def{} be.beautiful-\Pf{}//
\glft \trsl{The city became beautiful.} \lingcontext{perfective: ingressive}//
\endgl
\a
\begingl
\gla Nužo mrocnáv.//
\glb city.\Dir{}.\Def{} be.beautiful-\Ret{}//
\glft \trsl{The city was beautiful.} \lingcontext{retrospective: past state}//
\endgl
\a
\begingl
\gla Nužo mrocnáž.//
\glb city.\Dir{}.\Def{} be.beautiful-\Ctpl{}//
\glft \trsl{The city will be beautiful.} \lingcontext{contemplative: expected state}//
\endgl
\xe

In bounded or developmental aspects --- the progressive, perfective, and
retrospective --- stative verbs are normally interpreted ingressively unless
context strongly resists it. The perfective thus yields \trsl{became
beautiful}, not simply \trsl{was beautiful}; that reading belongs to the
retrospective, which describes a past state conceived as a complete episode, a
state that no longer obtains, or one that is no longer at issue in the current
discourse. The continuous, by contrast, gives plain property predication
without any implication of change, and is the unmarked form for simple
description.

The stative verbaliser \ird{-n-} selects the high-vowel perfective
\ird{-ik}; ordinary Class~A root G1 gives \ird{mrócnik} and
\ird{svórnik}. Its segmental suffix specification persists in
stative applicatives, although an intervening valency head prevents a
new aspectual root grade and supplies its own local quantity.
Selected quantities are retained before linker \ird{-e} and nominaliser
\ird{-ou} (§\ref{sec:aspect-pf}).

\subsection{Attributive and adverbial derivatives}\label{sec:stative-attr}

Unlike regular verbs, stative verbs derived from property roots possess a
dedicated attributive construction with lexically specified root-strong
and suffix-strong classes. Root-strong \ird{mrócna} selects G1 on the
property root and short \ird{-a}, while suffix-strong \ird{hradná},
\ird{svorná}, and \ird{skorná} retain their base grade with long
\ird{-á}. The attributive forms a separate derivational domain: protected
\ird{tvírzná} can retain both the lexical root quantity and long
attributive ending. These are class-selected correspondences, not a general
shortening of every ending after a long root. The construction is tenseless,
carries no aspectual or alignment marking, and is used exclusively as a
direct modifier within the noun phrase:

\pex
\a \ird{mrócna nuž} \trsl{a beautiful city}
\a \ird{hradná byl} \trsl{a good child}
\a \ird{svorná dum} \trsl{a big house}
\xe

The attributive form cannot be used predicatively. For predication, the full
stative verb form is required:

\pex
\a \ird{Nužo mrocna.} \trsl{The city is beautiful.} \lingcontext{predicative: stative verb required}
\a \ird{*Nužo mrócna.} \lingcontext{ungrammatical: attributive form in predicative position}
\xe

\noindent The attributive functions as the adnominal counterpart of
the continuous stative, expressing the property directly without tense or
aspect. Its position within the noun phrase relative to determiners and other
modifiers is treated in Chapter~\ref{ch:nouns}.

Alongside attributives in \ird{-a}/\ird{-á}, property-root-based stative stems
productively form a dedicated class of derived manner adverbs in \ird{-ě}
\label{sec:stative-adv}. This formation is derivational rather than
inflectional: it yields a lexical adverbial adjunct, not a further verbal
form, and it is distinct both from the adnominal attributive and from the
converb system discussed in Chapter~\ref{ch:complex}. Schematically, the
pattern is as follows:

\begin{center}
\term{property root $\rightarrow$ stative stem in \ird{-n-} $\rightarrow$ attributive in \ird{-a}/\ird{-á} / manner adverb in \ird{-ě}}
\end{center}

\noindent A representative paradigm is the following:

\pex
\a \ird{skor} \trsl{quickness}
\a \ird{skorn-} \trsl{be quick}
\a \ird{skorná} \trsl{quick, fast} \lingcontext{attributive}
\a \ird{skorňe} \trsl{quickly}
\xe

The resulting system preserves a clear three-way distinction: the finite stative
predicate expresses property predication, the attributive modifies a noun
within the \synhead{DP}, and the \ird{-ě} form modifies the event as a manner
adverb. Thus \ird{skorn-} means \trsl{be quick}, \ird{skorná} means \trsl{quick}
as a noun modifier, and \ird{skorňe} means \trsl{quickly} as a clausal adjunct.
The syntactic position of these derived manner adverbs in the preverbal field is
discussed in Chapter~\ref{ch:simple}; compare also the simultaneous / manner
converb in \ird{-ěc} in Chapter~\ref{ch:complex}, which remains a non-finite
verbal form rather than a lexical adverb.

\paragraph{Degree prefixes.}
Scalar degree on the stative class is expressed by a closed set of prefixes
that attach to the stative stem before any derivational material, including
the verbalising \ird{-n-} and attributive \ird{-a}/\ird{-á}. The five prefixes are
mutually exclusive and occupy a single slot:

\begin{table}[ht]
  \footnotesize
  \caption{Degree prefixes on stative verbs.}\label{tab:degpfx}
  \begin{tabularx}{\textwidth}{l l L L}
    \toprule
    {\sc prefix} & {\sc value} & {\sc finite} & {\sc attributive} \\
    \midrule
    \ird{po-}  & slight; somewhat      & \ird{pomrocna}  & \ird{pomrócna}  \\
    \ird{do-}  & moderate; quite       & \ird{domrocna}  & \ird{domrócna}  \\
    \ird{pre-} & strong; very          & \ird{premrocna} & \ird{premrócna} \\
    \ird{od-}  & elevated; exceedingly & \ird{odmrocna}  & \ird{odmrócna}  \\
    \ird{nad-} & excessive; too        & \ird{nadmrocna} & \ird{nadmrócna} \\
    \bottomrule
  \end{tabularx}
\end{table}

\noindent The prefixes are fully productive across the stative class. Because
they precede the stem complex, they appear before \ird{-n-} in finite forms
and before the root in attributives. Each is integrated into the lexical
prosodic word and bears its initial stress: \ird{premrocna} and
\ird{premrócna} are stressed on \ird{pre-}, irrespective of the long root
vowel in the root-strong attributive. With external negative \ird{za-}, the stressed
syllable of \ird{zapremrocna} remains \ird{pre-}
(§\,\ref{sec:stress}).

\paragraph{The emphatic suffix \ird{-ěsk-}.}
The suffix \ird{-ěsk-} selects the front grade, so a stative
\ird{-n-} appears as \ird{-ň-} before its vowel: \ird{mrocňeská}.
It is distinct from the degree prefix system and does not occupy the same
grammatical slot. It is restricted to the suffix-strong \ird{-á} attributive.
Its listed quantity is preserved in \ird{mrocňeská} and
\ird{premrocňeská}; attributive selection does not reach through
\ird{-ěsk-} to lengthen the root. The construction carries a marked, heightened, or literary flavour rather than a simple
scalar value --- closer in function to an emphatic particle than to a degree
adverb. The contrast with a prefixed form is one of register and framing
rather than degree alone: \ird{premrócna} states a degree as a plain fact;
\ird{mrocňeská} colours the attribution, implying that the beauty is
remarkable, conspicuous, or worth dwelling on. In ordinary speech
\ird{-ěsk-} attributives may sound affected; in verse, formal oratory, and
deliberate rhetorical prose they are appropriate and expected.

The suffix is not fully productive. It is most natural with roots that carry
inherent aesthetic or evaluative content:

\pex
\a \ird{mrocňeská} \trsl{strikingly beautiful; beautiful indeed} \lingcontext{natural in literary register}
\a \ird{hradňeská} \trsl{truly good; good in the full sense} \lingcontext{natural in literary register}
\xe

\noindent With dimensional or manner-denoting roots the formation is possible
but tends to sound stilted or unidiomatic outside verse:

\pex
\a \ird{svorňeská} \trsl{(strikingly big)} \lingcontext{odd in prose; possible in elevated style}
\a \ird{skorňeská} \trsl{(remarkably quick)} \lingcontext{marginal; feels forced outside verse}
\xe

\noindent \ird{trpěská} occupies a middle ground: the negative valence of
\ird{trpa} (\trsl{badness}) makes its \ird{-ěsk-} form feel ironic or
rhetorical rather than sincerely emphatic, and its use is consequently rare
and stylistically loaded.

\ird{-ěsk-} may co-occur with a degree prefix. The two are additive rather
than competing: the prefix supplies the scalar value and \ird{-ěsk-} supplies
the rhetorical register.

\pex
\a \ird{premrocňeská} \trsl{very beautiful} \lingcontext{emphatic; degree stated and dwelt upon}
\a \ird{odmrocňeská} \trsl{exceedingly beautiful} \lingcontext{elevated register; characteristic of panegyric}
\a \ird{nadmrocňeská} \trsl{excessively, ruinously beautiful} \lingcontext{ironic or elegiac; beauty to a fault}
\xe

\noindent Such combinations are licensed but marked; they are characteristic
of formal verse and deliberate oratorical prose. In ordinary speech, a
prefixed form without \ird{-ěsk-} is the normal choice.

\subsection{Nominalisation of stative verbs}

Stative predicates are nominalised by the same \ird{-ou} mechanism that
applies to other verbal predicates (§\,\ref{sec:nominalised}). A continuous
stative nominal denotes a participant characterised by the state, as in
\ird{mrocnova} \trsl{a beautiful one}. In the following equative, the
context supports identification with a relevant bearer of the property:

\pex
\a \ird{Nužo mrocna.} \trsl{The city is beautiful.} \lingcontext{simple property predication}
\a \ird{Nužo mrocnova.} \trsl{The city is the beautiful one.} \lingcontext{substantive: identifying or contrastive}
\xe

\pex
\a \ird{Bylo svórnik.} \trsl{The child became big (i.e.\ grew up).} \lingcontext{ingressive perfective}
\a \ird{Bylo svórnikou.} \trsl{The child is the one who grew up.} \lingcontext{nominalised perfective}
\xe

A non-continuous stative relative selects the same aspectual predicate as
the participant nominal, with relative \ird{-ový} in place of nominal
\ird{-ou}:

\ex
\begingl
\gla Svórnikový bylo Marek.//
\glb be.big-\Pf{}-\textsc{rel} child.\Dir{}.\Def{} Marek.\Dir{}//
\glft \trsl{Marek is the child who grew up.}//
\endgl
\xe

The continuous is an exception: the continuous stative nominalised in
\ird{-ova} cannot be used adnominally with \ird{-ový}. In adnominal position
the attributive construction must be used instead:

\pex
\a \ird{mrócna nuž} \trsl{a beautiful city} \lingcontext{correct: root-strong attributive}
\a \ird{*mrocnový nuž} \lingcontext{excluded continuous stative relative}
\xe

\subsection{State types and aspectual interpretation}\label{sec:stative-stability}

Not all states predicated by stative verbs are equally permanent. Some roots
denote enduring or characteristic properties---canonical evaluative qualities
such as \ird{hrad} \trsl{goodness} or \ird{mroc} \trsl{beauty}, and
dimensional properties such as \ird{svor} \trsl{bigness}---while others denote
contingent or temporary conditions: \ird{ptakn-} \trsl{be tired}, \ird{smrn-}
\trsl{be sick}, \ird{hvodn-} \trsl{be open}, \ird{žbirn-} \trsl{be wet}. This
distinction is not encoded grammatically, but it shapes the naturalness and
interpretation of aspectual forms.

The general contrast between continuous and progressive---laid out for all
predicate classes in §\,\ref{sec:aspect-ipf}---takes a particularly sharp form
with stative predicates, because here the interpretive difference between the
two is categorical rather than merely gradient. For a stable property, the
continuous yields plain predication and the progressive is distinctly marked,
carrying a developmental reading that implies a process of change
(\trsl{getting better}, \trsl{getting bigger}). For a temporary bodily or
experiential condition, the continuous describes the present condition as
simply obtaining (\trsl{is tired right now}), while the progressive more
readily receives an ingressive reading (\trsl{is starting to feel tired},
\trsl{is getting sick}). In both cases the continuous is the default
imperfective, giving simple present-state predication, while the progressive
imposes the dynamic, unfolding interpretation that is its characteristic
contribution across the predicate system.

Three semantic classes within the stative system may accordingly be
distinguished: \term{canonical quality statives}, which denote evaluative
properties of things (\trsl{be good}, \trsl{be beautiful}, \trsl{be bad});
\term{dimensional statives}, which denote extent or scale properties
(\trsl{be big}); and \term{experiential and bodily-condition statives}, which
denote states of persons (\trsl{be ill}, \trsl{be hungry}, \trsl{be ashamed}).
The last class most readily admits progressive and retrospective readings, since
temporary and personally felt states support aspectual dynamics more naturally
than stable properties do.

\subsection{Valency operations with stative verbs}\label{sec:stative-valency}

Stative causatives describe the bringing about of a state or a change along
the scale expressed by the base. The progressive presents that change as
unfolding, whereas the perfective presents a bounded causing event; neither
viewpoint specifies whether the causer acts directly or through an intermediary.
Stative stems in \ird{-n-} select expanded causative \ird{-er-}, as in
\ird{svorn-ér-ime}. Overt outer Voice independently takes \ird{-an-}
after the causative, giving perfective \ird{svorn-ér-an-ik} and
retrospective \ird{svorn-ér-an-av}. The stem-final stative \ird{-n-},
the rhotic causative, and outer orientation therefore occupy distinct
positions, without metathesis or a joint causative--Voice exponent.

\pex\label{ex:val-stative-caus}
\a
\begingl
\gla Marek problemam svornérime.//
\glb Marek.\Dir{} problem.\Ind{} be.big-\Caus{}-\Prog{}//
\glft \trsl{Marek is making the problem bigger.}//
\endgl
\a
\begingl
\gla Problem Marku svornéranik.//
\glb problem.\Dir{} Marek.\Trs{} be.big-\Caus{}-\textsc{trs}-\Pf{}//
\glft \trsl{Marek made the problem bigger.}//
\endgl
\xe

Constructed continuous \ird{svornéra} \trsl{makes bigger} and
retrospective \ird{svornéranav} \trsl{has made bigger} follow the same
NOM/ERG case contrast. Their forms preserve the same expanded causative and the ordinary
conditioning of outer Voice.

Stative applicatives require a lexically licensed relation to the state.
The stimulus construction with \ird{všihn-} \trsl{be happy} relates an
experiencer to the news about which the state holds. Its applied stimulus is
direct and its experiencer transitive in the perfective:

\ex
\begingl
\gla Senica Marku všihnébik.//
\glb news.\Dir{} Marek.\Trs{} be.happy-\Appl{}-\Pf{}//
\glft \trsl{Marek became happy about the news.}
\lingcontext{lexically restricted stimulus applicative}//
\endgl
\xe

The perfective presents entry into the state, while the applied stimulus
identifies what Marek is happy about. Neither this stimulus relation nor
perfective aspect asserts that the news caused the state to arise; a causal
relation requires causative morphology or an independently established lexical
causative.
The stimulus extension supplies no general rule admitting a causer with every
stative applicative.

Underived stative verbs are intransitive and lack the lexical patient required
by the antipassive. Their derived causatives and applicatives license additional
participants through their respective argument-building heads; those
derivations do not establish a patient-bearing reading of the stative base.

Stative verbs may further function as attributive modifiers of nouns and remain
the productive attributive strategy in Iridian. Alongside them stands the small
closed set of paired adnominal property-root forms just mentioned; this use is
described in Chapter~\ref{ch:nouns}.

Property-root-based stative stems also feed the productive derived manner
adverbs in \ird{-ě} (§\,\ref{sec:stative-adv}), yielding lexical adverbial
adjuncts such as \ird{skorňe} \trsl{quickly}. These belong to the clause's
preverbal adjunct structure, not to the attributive system proper.

\subsection{Syntactic analysis of stative verbs}
\label{sec:vm-tension-stative}

Stative verbs inflect for the same aspectual distinctions as other finite
verbs. A lexical property root combines
with the verbalising suffix \ird{-n-}, analysed syntactically as a stativising
verbal head, to yield an intransitive predicate denoting a state. That
predicate is then selected by the same Asp head that selects other finite
verbs in Iridian, so the full five-way Asp.Align system applies to stative
bases exactly as it does elsewhere in the verbal system.

In the continuous, the stative predicate denotes a state holding at the
reference time, as in \ird{Nužo mrocna}
\trsl{the city is beautiful}. In the progressive, perfective, retrospective,
and contemplative, the same Asp system contributes developmental, bounded, or
temporally displaced viewpoints over that stative base; the resulting readings
are therefore naturally inchoative or ingressive (\trsl{be becoming
beautiful}, \trsl{became beautiful}) rather than simple present-state
descriptions, with the aspectual contribution interpreted over the lexical
state predicate.

The distinction between a finite stative predicate and a nominal predicate is
also reflected in negation, even where neither clause contains an overt copula:

\begin{center}
\begin{tabular}{l l l}
\toprule
\textsc{predication type} & \textsc{affirmative} & \textsc{negative} \\
\midrule
Zero-copula equative & \ird{Marek doktor} & \ird{Marek doktor staly} \\
Stative verbal & \ird{Nužo mrocna} & \ird{Nužo zamrocna} \\
\bottomrule
\end{tabular}
\end{center}

\medskip
\noindent The nominal predicate \ird{doktor} takes clause-final
\ird{staly}, whereas \ird{mrocna}, formed from \ird{mrocn-} with continuous
\ird{-a}, takes \ird{za-}. Both clauses lack an overt copula, but the stative
predicate bears verbal and aspectual morphology.

Neg is analysed as selecting a nominal predication in the equative and an
aspectually inflected verbal predicate in the stative clause
(§\,\ref{sec:vm-neg}). Finite property predicates contain the verbalising
suffix \ird{-n-} and inflect for Asp.Align; predicative use of underived
property forms is undocumented.

The syntactic infrastructure of the stative domain therefore distinguishes
three outputs from the same lexical base. First, the finite stative verb heads
the clause's verbal predicate and requires no copula in ordinary property
predication. Second, the \ird{-á} form is a non-finite attributive derivative
inside the \synhead{DP}, corresponding to the adnominal strategy discussed in
Chapter~\ref{ch:nouns}; it is not a finite predicate and cannot bear aspect or
alignment. Third, the \ird{-ě} form is a lexical adverbial derivative that
functions as a clausal adjunct in the preverbal field, not as an inflectional
member of the verbal paradigm. The same lexical base thus feeds clause-level
predication, \synhead{DP}-internal attribution, and adverbial modification through three
distinct syntactic realisations.

Because underived stative predicates are intransitive, the alignment component
of Asp.Align is trivial in their simple finite use: the sole argument is the
direct-case absolutive/nominative argument in every aspect. The ordinary
alignment alternations re-enter the system only when valency-changing
morphology is added. Causativised statives behave like other derived
transitives, showing nominative alignment in the progressive and contemplative
and ergative alignment in the perfective and retrospective. The stative system
thus does not challenge the broader verbal theory; it instantiates it with a
lexically stative base whose aspectual interpretations follow from the
interaction of a regular Asp head with a regular intransitive predicate.


\section{Verbal complements and head finality: preverbal position}
\label{sec:vm-comp-headfinal}

This section develops the formal structural analysis of verbal complements within the head-final clause: how the head-final parameter determines the preverbal position of nominal arguments and embedded clauses, and what syntactic projection each complement type occupies.

\subsection{The base position of the object}
\label{sec:vm-obj-base}

In a head-final language, the complement of a head $H$ is to the
left of $H$ in the linear string. For the lexical predicate V, its
internal argument (patient \synhead{DP}) is the complement of V and therefore
surfaces to V's left:

\begin{center}
\term{[\synhead{VP} [\synhead{DP} patient] V]}
\end{center}

\noindent This is a \emph{base-generated} preverbal position, not
the result of A-bar movement or scrambling. In the Minimalist
framework, head-finality is parameterised as the linearization
convention that assigns $<$Complement, Head$>$ ordering (Kayne's
1994 antisymmetric alternative --- deriving all
SOV orders from movement --- is not adopted here; the parametric
head-final approach is more consistent with the thoroughgoing
head-final organisation of Iridian lexical phrases).

The patient \synhead{DP} is base-generated as the complement of V, before
 the verb. Its direct or indirect case is licensed within the aspectual and
 verbal domain; case assignment does not require movement to
 Spec,\synhead{TopP}. A lower initial pivot can consequently lack an
 aboutness interpretation (§\,\ref{sec:tc-quantificational-pivots}).

When an argument is independently selected as a linked topic, an
 A$'$-dependency connects Spec,\synhead{TopP} with that lower position.
 Schematically, with intervening projections suppressed:
\[
 [_{\mathrm{TopP}}\ \mathrm{DP}_{i}\
   [\mathrm{lower\ clause}\ \ldots t_i\ldots V]\ \mathrm{Top}^{0}].
\]
The \synhead{DP} retains its case, so a perfective transitive can have either
 a direct-case patient topic or a transitive-case agent topic. Head-finality
 places complements before V; the topic dependency, when present, places its
 referent above the lower clause. These are distinct sources of preverbal
 order.

\subsection{The po-clause as a nominal \synhead{CP} complement}
\label{sec:vm-po-comp}

The finite propositional complement (\ird{po}-clause) is a \synhead{CP} that
occupies the complement position of the matrix predicate. The prototypical
structure of the complement \synhead{CP} is:

\begin{center}
\term{[\synhead{CP} [\synhead{TP} embedded-clause-\ird{e}] \ird{po}]}
\end{center}

\noindent \ird{Po} is the head of the complement \synhead{CP} (a
complementizer), and the embedded finite clause is the \synhead{TP} whose
head is marked by the linker \ird{-e}. In head-final structure,
the C head follows its complement \synhead{TP}:

\begin{center}
\begin{forest}
  syntax tree
  [\synhead{CP}
    [{\synhead{TP}\\\textit{embedded clause}\\\textit{(verb + -e)}}]
    [{C\\\ird{po}}]
  ]
\end{forest}
\end{center}

\noindent This \synhead{CP} is then the complement of the matrix verb:

\begin{center}
\begin{forest}
  syntax tree
  [\synhead{VP}
    [{\synhead{CP}\\{[\textit{emb-clause}-\ird{e}] [\ird{po}]}}]
    [V]
  ]
\end{forest}
\end{center}

\noindent Head-finality places the \synhead{CP} complement to the left of the
matrix V. The surface word order of a propositional complement
construction is therefore:

\begin{center}
\term{[embedded-\synhead{TP} + \ird{-e}] [\ird{po}] [matrix arguments] [matrix V]}
\end{center}

\noindent This is the order illustrated (§\,\ref{sec:complementation}).
The analysis requires no movement: the complement \synhead{CP} is
base-generated in [Complement, matrix-\synhead{VP}], and head-finality places
it before the matrix V.

When the embedded predicate is a zero-copula predication, however, no overt
finite verb is available to host the boundary suffix \ird{-e}. In that special
case Iridian inserts the support form \ird{ce} immediately before \ird{po}
(and, in the corresponding non-assertive construction, before \ird{poby}).
This \ird{ce} is best analysed not as a general copular verb but as a
clause-support exponent used to realize the complement edge for otherwise
verbless finite predications. The resulting pattern is therefore:

\begin{center}
\term{[\synhead{CP} [small clause + \ird{ce}] \ird{po}]}
\end{center}

A recoverable matrix argument may remain unexpressed when the discourse
 identifies its referent. The propositional complement and an overt matrix
 topic occupy distinct positions, so the presence or weight of the
 \synhead{CP} does not syntactically exclude an overt topic. Null realisation
 follows the general recoverability conditions
 (§\,\ref{sec:tc-topic-drop}), independently of complement placement.

\paragraph{Po as C$^0$: evidence.}

The analysis of \ird{po} as a C head (and not, for example, as an
adposition or particle) is supported by the following diagnostics:
\begin{enumerate}
\item The embedded clause is fully finite (retains aspect,
  alignment and mood morphology); this is what is
  expected of a \synhead{CP} complement, not a reduced \synhead{VP}.
\item \ird{Po} appears after the embedded clause-final verb
  (the right edge of \synhead{TP} in a head-final \synhead{CP}), in exactly the
  position predicted for a head-final C.
\item \ird{Poby} (the non-assertive counterpart) parallels \ird{po}
  structurally and has the semantics of a complementizer: it
  types the embedded clause as an indirect interrogative
  (\textit{whether} or embedded wh-question) rather than an
  assertive declarative. Complementizer alternation is a
  \synhead{CP}-level phenomenon.
\item The linker \ird{-e} on the embedded verb is the edge-marking
  morpheme at the \synhead{TP}-\synhead{CP} boundary; where no verbal host is
  present because the predicate is zero-copular, the support form \ird{ce}
  appears instead. This is consistent with analysing the complement edge as a
  suffixal boundary marker on finite verbal clauses and as a support position
  on finite verbless predications (§\,\ref{sec:vm-tension-linker}).
\end{enumerate}

\subsection{The purposive converb complement}
\label{sec:vm-purp-comp}

The purposive converb (\ird{-it}) serves as the complement of
desiderative, volitional, modal and motion predicates. Unlike the
\ird{po}-clause, it lacks finite viewpoint and mood, although its internal
arguments receive the NOM-domain cases licensed by the converb construction.
Its lower predicate can contain applicative and causative morphology; the
purposive head Conv, realised by \ird{-it}, closes this constituent without
introducing finite Asp.Align (§\,\ref{sec:vm-converb-aspless}).

The purposive complement is analysed as \synhead{ConvP}, whose head selects
the lower \synhead{vP}/\synhead{VoiceP}, rather than as a finite
\synhead{CP}. It denotes a prospective event description, while a
\ird{po}-complement presents a proposition. In head-final structure, Conv
follows its verbal complement and the entire purposive constituent precedes
the matrix predicate:

\begin{center}
\term{[matrix pivot] [\synhead{ConvP} [lower verbal predicate] \ird{-it}] [matrix V]}
\end{center}

\noindent The implicit participant of the purposive complement normally
shares the reference of the matrix absolutive pivot
(§\,\ref{sec:switch-reference}). In the ordinary intransitive-motion
examples, that pivot is the mover, so the dependency yields a same-subject
reading (§\,\ref{sec:purp-motion}). This reference relation is independent
of aboutness topichood and does not require an additional phonologically null
finite Asp head.


\subsection{The preverbal zone: a structural hierarchy}
\label{sec:vm-preverbal-zone}

Within the matrix clause, the preceding analysis distinguishes the following
fields in their surface order. Selected clauses and converbs have
construction-specific positions, described separately below the table.

\begin{center}
\footnotesize
\begin{tabular}{l l l}
  \toprule
  {\sc position} & {\sc constituent type} & {\sc structural slot} \\
  \midrule
  1 (leftmost) & Zone I peripheral particles & Spec of \synhead{EvidP}, \synhead{SpeechActP} \\
  2            & Topic \synhead{DP} (case retained)     & Spec,\synhead{TopP} \\
  3            & Clause-typing particle \ird{by} & Head of \synhead{PolP} \\
  4            & Wh-phrase (in questions)  & Spec,\synhead{FocP} \\
  5            & Non-topic arguments       & Spec,\synhead{IP} or \synhead{IP}-adjuncts \\
  6            & Comment-domain operators   & Spec/Head of \synhead{AspP} (zone II) \\
  7            & Phasal particles           & Spec/Head of \synhead{PhasP} (zone III) \\
  8 (rightmost) & Finite verb               & Head of \synhead{VP} / \synhead{vP} \\
  \bottomrule
\end{tabular}
\end{center}

\noindent A selected clause or converb belongs to the lexical predicate's
complement domain, but its base position does not establish a single surface
slot shared by all complement constructions. Finite \ird{po}-complements
precede the matrix arguments in the pattern illustrated above; purposive
constructions can place the matrix topic before the converb
(§\,\ref{sec:purp-motion}). Neither pattern requires the complement to be
immediately adjacent to the finite matrix verb. Where a matrix phasal particle
is present, the ordinary predicate-adjacency rule predicts that the complement
precedes that particle.



%% ================================================================
%%  SUMMARY TABLE
%% ================================================================

\section{Summary: the verbal complex as a syntactic structure}
\label{sec:vm-summary}

Table~\ref{tab:vm-full-structure} gives the hierarchical structure
of the Iridian finite verbal clause, from the highest projection
(\synhead{SpeechActP}) to the lowest (\synhead{VP}/Root), with the morphological
exponents of each relevant layer.

\begin{table}[ht]
  \footnotesize
  \caption{The Iridian verbal clause: functional hierarchy and
  morphological exponents.}\label{tab:vm-full-structure}
  \begin{tabularx}{\textwidth}{l l l L}
    \toprule
    {\sc level} & {\sc projection} & {\sc morphological exponent}
    & {\sc notes} \\
    \midrule
    Highest  & \synhead{SpeechActP} & Zone I particles (peripheral)
    & Upper-peripheral speech-act framing \\
    & \synhead{EvidP}$_\text{outer}$ & \ird{oče, hlavdy, ktady, či}
    & Clause-level evidential/epistemic \\
    & \synhead{TopP} & [topic \synhead{DP}]
    & Aboutness topic; case retained from lower domain \\
    & \synhead{PolP} & \ird{by}
    & Head-initial clause typing: interrogative \\
    & \synhead{FocP} & [wh-phrase]
    & Focus/wh movement target \\
    \midrule
    & \synhead{EvidP}$_\text{inner}$ & \ird{-al}, postvocalic \ird{-l}
    & Indirect evidential; inside verbal complex \\
    & \synhead{MoodP} & \ird{-ím, -é, -ká, -ku, -ěv, -ir}, $\varnothing$
    & Negative Mood bears $[u\textsc{neg}]$; reduces Asp features \\
    & \synhead{NegP} & \ird{za-} (prefix)
    & Op$_{\neg}$ in specifier; Neg$^{0}[u\textsc{neg}]$; PF attachment \\
    & \synhead{AspP} & \ird{-ek/-ik, -av, -a, -ime, -áž} (fused w/ align)
    & Sole obligatory head; aspect + case assignment \\
    \midrule
    & \synhead{VoiceP}$_{\mathrm{orient}}$ & \ird{-n-}/\ird{-an-}, $\varnothing$
    & Selects participant of the complete derived predicate \\
    & \synhead{CausP} & \ird{-r-}, \ird{-er-}, \ird{-ěr-}
    & Adds causer above the inner valency domain \\
    & \synhead{VoiceP} / \synhead{ApplP} & \ird{-aš-} (antipassive) / \ird{-éb-} (applicative)
    & Distinct heads; language-specific co-occurrence ban \\
    & \synhead{VP} & Root stem
    & Lexical predicate; complement of V = preverbal \\
    Lowest & \synhead{VP} complement & \synhead{NP} / \synhead{CP} (\ird{po}) / \synhead{VoiceP}$_{\text{-it}}$
    & Object, po-clause, or purposive converb; all preverbal \\
    \bottomrule
  \end{tabularx}
\end{table}

\noindent The preverbal position of all verbal complements --- \synhead{NP}
objects, \synhead{CP} complements (\ird{po}-clauses), and purposive converb
complements --- follows uniformly from the head-final parameter applied at the
\synhead{VP} level. No movement or extraposition is required: the complement is
base-generated to the left of the verb and surfaces there in the linear string.
Asp, Mood and the inner evidential have suffixal exponents on the verbal
complex, which occupies the right edge of the lower clause. Neg is structurally
higher than Asp but attaches at the left edge of the verbal host, as specified
in §\,\ref{sec:vm-neg}.

%% ----------------------------------------------------------------
%% END: verb-minim.tex
%% ----------------------------------------------------------------
